\chapter{Experimental Supplement}
\label{sec:supplement}

For each chapter, this appendix collects its main claims and what they rest on: the experiment in which each was measured, what exactly was measured, and where it was published. The status in the right-hand column says how firm a claim is: \emph{measured} means there is a direct experiment; \emph{theory}, a mathematical consequence derived in the chapter or in a classic paper; \emph{model}, a result computed with the book's model (scripts in the \texttt{models/} folder); \emph{estimate}, an order-of-magnitude number without direct measurement; \emph{hypothesis}, a plausible explanation not yet proven by experiment. Where there is no experiment, we say so.

\section{Chapter~\ref{sec:neurontypes}. Neuron Types}

The chapter's numbers rest on direct cell counts in the human brain where such counts exist; the ``one neuron, one transmitter'' rule rests on cell atlases and on experiments that have also found the exceptions; the role of \nt{DOPA} rests on spike recordings; and the roles of the other broadcast channels rest on hypotheses.

{\footnotesize
\setlength{\tabcolsep}{3pt}\renewcommand{\arraystretch}{1.15}
\begin{longtable}{>{\raggedright}p{0.5cm}>{\raggedright}p{4.0cm}>{\raggedright}p{5.6cm}>{\raggedright}p{2.3cm}>{\raggedright\arraybackslash}p{1.7cm}}
\toprule
No. & Claim & Experiment and what was measured & Source & Status \\
\midrule\endfirsthead
\toprule
No. & Claim & Experiment and what was measured & Source & Status \\
\midrule\endhead
1 & The human brain has about $8.6\cdot10^{10}$ neurons, and almost all \nt{GLUT} neurons are small cerebellar neurons (Table~\ref{tab:types}) & Isotropic fractionator: tissue is dissolved into a uniform suspension of nuclei, and nuclei carrying the neuronal marker NeuN are counted. Adult men have $86.1\pm8.1$ billion neurons; only $19\,\%$ of them are in the cerebral cortex, and the bulk is in the cerebellum & \cite{Azevedo2009} & measured \\
2 & Almost every neuron has one principal transmitter (Proposition~\ref{prop:dale}), but there are exceptions & Transcriptomic atlas of the mouse brain: about $4\cdot10^6$ cells, $5322$ clusters; each cluster is assigned a transmitter from its synthesis and packaging genes. The exceptions were measured directly: in slices, \nt{DOPA} neurons also release GABA through the same vesicular transporter and inhibit striatal neurons; in some \nt{DOPA} axons, glutamate and dopamine leave from different terminals of the same wire (electron microscopy and optogenetics) & \cite{Strata1999,Yao2023,Tritsch2012,Zhang2015} & measured \\
3 & The entire column of the weight matrix has one sign~\eqref{eq:dalesign} & Derived in the chapter from Proposition~\ref{prop:dale} and the fact that glutamate receptors are almost all excitatory and GABA receptors inhibitory. There is no direct test of ``all recipients of one neuron receive a weight of one sign'' as a general rule; counterexamples are GABA co-release by \nt{DOPA} neurons (row~2) and recipients with receptors of opposite sign (row~11) & derivation in the chapter & theory \\
4 & Three-quarters to four-fifths of cortical neurons are \nt{GLUT}, a fifth to a quarter are \nt{GABA} & GABA immunostaining in columns $50\um$ wide through the full depth of the cortex in $10$ areas of monkey cortex: in $8$ areas GABA neurons make up about $25\,\%$ of all neurons, in primary visual cortex less than $20\,\%$ & \cite{Hendry1987} & measured \\
5 & A \nt{GLUT} neuron has about $10^4$ outputs & Postmortem stereology: the neocortex of young men has $1.64\cdot10^{14}$ synapses (electron microscopy, disector) and about $2.3\cdot10^{10}$ neurons. The ratio gives $\sim7\cdot10^3$ synapses per neuron; on average the number of inputs equals the number of outputs & \cite{Tang2001synapses,Pakkenberg1997} & measured \\
6 & The output of a \nt{DOPA} neuron branches into $10^5$--$10^6$ terminals; this is an estimate from target coverage, not a count & Viral membrane labeling of single rat \nt{DOPA} neurons and full axon reconstruction: one neuron covers $0.45$ to $5.7\,\%$ (on average $2.7\,\%$) of the striatal volume. Coverage was measured, not the number of terminals: that is inferred from coverage to order of magnitude, with no direct terminal count. For \nt{SERO} and \nt{NORA} the terminal numbers are rough estimates & \cite{Matsuda2009} & measured \\
7 & Broadcast neurons are few: \nt{DOPA} $\sim5\cdot10^5$ (the main one of two groups, a lower bound), \nt{SERO} $\sim3\cdot10^5$ (sum of two partial counts), \nt{HIST} $\sim6\cdot10^4$ (Table~\ref{tab:types}, Fig.~\ref{fig:typepie}) & Stereological counts in humans: $550\,000$ pigmented neurons in the substantia nigra (without the ventral tegmental area); $235\,000$ neurons in the dorsal raphe nucleus (all neurons of the nucleus) and another $125\,000$ serotonin neurons in the pontine tegmentum; about $64\,000$ histamine neurons in the hypothalamus. For \nt{NORA}, \nt{GLYC}, \nt{ACET}, and \nt{ADRE} there is no direct count: in the table these are rough order-of-magnitude estimates & \cite{Pakkenberg1991,Baker1990,Baker1991,Airaksinen1991} & measured \\
8 & Glutamate in the gap shoots up to about a millimolar and decays within $\sim1\ms$ & From the displacement of a rapidly unbinding competitive blocker from NMDA receptors during synaptic transmission (cultured hippocampal neurons): peak $1.1$\,mM, decay time constant $1.2\ms$ & \cite{Clements1992} & measured \\
9 & In working cortex the excitatory and inhibitory currents are nearly balanced, and the neuron sits near threshold & Theory: a network with strong connections settles into the balanced regime on its own. Experiment: in ferret prefrontal cortex in vivo, during ``up'' states the reversal potential of the synaptic current stays near $-37$\,mV for hundreds of milliseconds, although the conductance changes by $\sim21$\,nS: excitation and inhibition rise and fall together & \cite{vanVreeswijk1996,Haider2006} & measured \\
10 & \nt{DOPA} neurons report the reward-prediction error~\eqref{eq:rpe} (Fig.~\ref{fig:rpe}) & Spikes of monkey \nt{DOPA} neurons: a burst for unexpected juice; after learning, a burst for the cue and no response to predicted juice; a dip in rate when promised juice is withheld. In a quantitative experiment, the rate is proportional to the difference between the current reward and a weighted average of past ones, but only for positive errors. Equation~\eqref{eq:rpe} itself is the temporal-difference algorithm & \cite{Schultz1997,Bayer2005,SuttonBarto2018} & measured \\
11 & The receptor sets the sign and speed: ionotropic, fractions of a millisecond; metabotropic, much slower (Proposition~\ref{prop:receptor}) & Brief pulses of glutamate on membrane patches from hippocampal neurons: the current through ionotropic receptors rises (from $20$ to $80\,\%$) in $0.2$--$0.6\ms$ and decays in $2$--$3\ms$. The slow inhibitory potential of pyramidal neurons is abolished by a selective blocker of the metabotropic GABA receptor. Two families of dopamine receptors with opposite action; in the striatum, neurons with the D1 receptor need dopamine to strengthen synapses, neurons with D2 to weaken them & \cite{Colquhoun1992,DutarNicoll1988,Beaulieu2011,Shen2008} & measured \\
12 & A broadcast channel is not one wire but a bundle of a few lines (Section~\ref{sec:buses}) & Viral-genetic labeling of serotonin neurons in the mouse dorsal raphe nucleus: neurons projecting to the amygdala and to the frontal cortex lie in different parts of the nucleus, branch into complementary regions, and respond oppositely to an aversive stimulus. Review: subgroups of locus coeruleus neurons (\nt{NORA}) encode different processes & \cite{Ren2018,Poe2020} & measured \\
13 & \nt{NORA} sets the gain $g$ & The adaptive gain hypothesis; a network model in which catecholamines increase the steepness of the transfer characteristic. There is no direct measurement of ``$g$ rises with noradrenaline'' across the whole network; what has been measured is that pupil diameter follows the spikes of locus coeruleus neurons in the monkey & \cite{AstonJones2005,ServanSchreiber1990,Joshi2016} & hypothesis \\
14 & \nt{ACET} switches ``write/read'' and sets the learning rate & Hypothesis. Support: in rat olfactory cortex slices, acetylcholine agonists ($100$\,\textmu M) weaken the internal, ``recalling'' synapses by more than $60\,\%$ and the input synapses by less than $15\,\%$ & \cite{Hasselmo2006,HasselmoBower1992} & hypothesis \\
15 & \nt{SERO} sets the discount factor $\gamma$; serotonin neurons fire steadily, a few spikes per second, and are almost silent in one of the phases of sleep & The ``serotonin $=\gamma$'' hypothesis. Strongest argument: optogenetic activation of serotonin neurons in the mouse dorsal raphe nucleus reduces the number of premature give-ups while waiting for a delayed reward. The firing rate and silence during rapid-eye-movement sleep have been measured (review of recordings) & \cite{Doya2002,Miyazaki2014,JacobsAzmitia1992} & hypothesis \\
\bottomrule
\end{longtable}}

\section{Chapter~\ref{sec:glutcomputer}. What \nt{GLUT} Neurons Compute}

The chapter's logical claims are theorems about circuits with nonnegative weights, derived in the chapter itself; experiment supports the neuron model and the fact that circuits built according to these theorems (the coincidence detector, the delay line, the self-sustaining group, the chain) really do occur in the brain.

{\footnotesize
\setlength{\tabcolsep}{3pt}\renewcommand{\arraystretch}{1.15}
\begin{longtable}{>{\raggedright}p{0.5cm}>{\raggedright}p{4.0cm}>{\raggedright}p{5.6cm}>{\raggedright}p{2.3cm}>{\raggedright\arraybackslash}p{1.7cm}}
\toprule
No. & Claim & Experiment and what was measured & Source & Status \\
\midrule\endfirsthead
\toprule
No. & Claim & Experiment and what was measured & Source & Status \\
\midrule\endhead
1 & A neuron is an RC circuit with a threshold and reset~\eqref{eq:glutneuron} & A noisy current resembling the synaptic bombardment in the living brain was injected into the somata of rat cortical pyramidal neurons (slices). The dependence of the mean rate on the mean and variance of the current is described by a leaky integrate-and-fire neuron with added slow rate adaptation. The ranges of $\tau$ and delays are given in the chapter without reference to an experiment & \cite{Rauch2003} & measured \\
2 & Model~\eqref{eq:glutneuron} is a simplification: the input branches generate local spikes of their own & Recording directly from dendrites of pyramidal neurons in mouse visual cortex in vivo: a visual stimulus evokes local dendritic spikes that depend on NMDA receptors; suppressing them broadens the neuron's orientation tuning & \cite{Smith2013} & measured \\
3 & The coincidence window is $\Delta_{\max}=\tau\ln\frac{w}{\theta-w}$~\eqref{eq:window}; for $\theta=1.5w$ it equals $0.7\tau$ & A consequence of model~\eqref{eq:glutneuron}. A direct measurement of a narrow summation window exists for neurons with fast inhibition (Chapter~\ref{sec:gaba}, row~3 of its table) & derivation in the chapter & theory \\
4 & A network of \nt{GLUT} neurons alone computes only monotone functions of its inputs (Proposition~\ref{prop:monotone}) & Proof in the chapter: iteration from silence with a nondecreasing right-hand side. This is a claim about the model; there is no experiment testing it on a living network without inhibition & derivation in the chapter & theory \\
5 & Sense organs provide a pair of wires: some cells respond to an increase of the stimulus, others to a decrease & Retina: already at the bipolar cells the cone signal splits into on and off channels. Pharmacological block of on bipolar cells in the monkey: the channels run separately up to the cortex; after the block, detection of light increments is impaired but not of decrements & \cite{Schiller1992} & measured \\
6 & With a dual-rail code, a network of \nt{GLUT} neurons computes any logical function asynchronously, without a common beat, at the price of twice as many neurons (Proposition~\ref{prop:dualrail}, \eqref{eq:dualrail}, Fig.~\ref{fig:dualxor}) & The dual-rail code and completion detection from the signal itself are a standard technique of delay-insensitive asynchronous circuits. The six-neuron exclusive-OR circuit is built in the chapter. There is no experiment showing that the brain computes logical functions specifically with a dual-rail code & \cite{Sparso2001} & theory \\
7 & The largest interaural arrival-time difference in humans is $\approx0.6\ms$, and the brain resolves about ten microseconds & Listeners in a two-alternative experiment: the time-difference discrimination threshold (75\,\% correct) is $9$\,\textmu s for noise in the $150$--$1700$\,Hz band, $11$\,\textmu s for a $1$\,kHz tone, $28$\,\textmu s for a click. The $0.6\ms$ limit is the chapter's calculation, $0.22\,\text{m}/343\,\text{m/s}$ & \cite{KlumppEady1956} & measured \\
8 & In birds, a time difference is turned into a place by delay lines and coincidence detectors~\eqref{eq:itd} (Fig.~\ref{fig:itd}) & Barn owl: axons from the two ears enter the nucleus of coincidence detectors from opposite sides; the arrival time of spikes along the axon changes systematically with depth, and the spread of delays across the nucleus ($160$\,\textmu s over $700\um$) matches the range of interaural delays & \cite{Carr1990} & measured \\
9 & In mammals no row of delays has been found; the same task is solved by coincidence detectors with precisely timed inhibition from \nt{GLYC} neurons & In vivo recording in the gerbil medial superior olive: the responses do not form a map of directions; iontophoresis of glycine and its blocker shows that precisely timed \emph{glycinergic} inhibition sets the working range of delays. The inhibition here comes from \nt{GLYC}, not from \nt{GABA} & \cite{Brand2002,Grothe2010} & measured \\
10 & A group with strong feedback is a flip-flop; self-sustaining activity persists for seconds after the stimulus (Fig.~\ref{fig:bistable}) & Monkey prefrontal cortex in a delayed eye-movement task toward a remembered location (delay $1$--$6$\,s): $87$ of $170$ task-related neurons change their rate during the delay, and in $79\,\%$ of them this depends on the direction of the remembered location. That this activity is held by feedback with two stable states is theory & \cite{Funahashi1989,Wang2001} & measured \\
11 & A bit is written if the drive lasts longer than $\approx0.7\tau$ (Fig.~\ref{fig:recurrentgroup}b) & Euler integration of $\tau\,\dd r/\dd t=-r+f(wr+I)$ with $w=1$, $I=0.6$; there is no script in \texttt{models/}. No experiment & calculation in the chapter & model \\
12 & Memory can be stored in synaptic facilitation, with almost no spikes & Theory: residual calcium in the terminals holds the record for about a second. Support: in ferret medial prefrontal cortex (simultaneous recording from several neurons) there is a subnetwork of pyramidal neurons connected by facilitating synapses with long aftereffects. A review of observations of ``quiet'' intervals during memory maintenance. Which way the cortex uses when is an open question & \cite{Mongillo2008,WangMarkram2006,Stokes2015} & hypothesis \\
13 & Memory is erased by fatigue: the stock of transmitter is depleted & Paired recordings of cortical pyramidal neurons: the synaptic response falls with frequent spikes, and the rate of the fall is set by the release probability. That this is precisely what turns off a self-sustaining group has not been shown directly & \cite{Tsodyks1997} & measured \\
14 & A chain of groups is an event-triggered timer; in songbirds neurons fire strictly in turn & Recording of identified neurons of the HVC song nucleus in the zebra finch during sleep and singing: each neuron projecting to the motor nucleus fires one short burst at one precise moment of the song, and the neurons fire in sequence & \cite{Hahnloser2002} & measured \\
\bottomrule
\end{longtable}}

\section{Chapter~\ref{sec:gaba}. \nt{GLUT} and \nt{GABA} Together}

The chapter's logical and dynamical claims are derived in the chapter itself from linear analysis and Ohm's law; experiment shows that the circuits they rely on (the delayed veto, direct inhibition following excitation, stabilization by inhibition, the rhythm of the \nt{GLUT}--\nt{GABA} loop) really do operate in the brain.

{\footnotesize
\setlength{\tabcolsep}{3pt}\renewcommand{\arraystretch}{1.15}
\begin{longtable}{>{\raggedright}p{0.5cm}>{\raggedright}p{4.0cm}>{\raggedright}p{5.6cm}>{\raggedright}p{2.3cm}>{\raggedright\arraybackslash}p{1.7cm}}
\toprule
No. & Claim & Experiment and what was measured & Source & Status \\
\midrule\endfirsthead
\toprule
No. & Claim & Experiment and what was measured & Source & Status \\
\midrule\endhead
1 & A fifth to a quarter of cortical neurons are \nt{GABA} & GABA immunostaining in $10$ areas of monkey cortex: about $25\,\%$ of neurons in $8$ areas, less than $20\,\%$ in primary visual cortex. A review of the diversity of cortical inhibitory neurons & \cite{Hendry1987,Markram2004} & measured \\
2 & What inhibition does is largely decided by its landing site: dendrite, cell body, spike initiation site, or the terminal of another neuron's wire & Review: classes of cortical \nt{GABA} neurons differ in their target on the recipient. Simultaneous recording from the soma and dendrite of a pyramidal neuron: feedforward inhibition is much stronger at the soma than at the dendrites. Inhibition at the terminal of another neuron's wire in the spinal cord reduces transmitter release from one input line (review of measurements) & \cite{Markram2004,Pouille2001,Rudomin1999} & measured \\
3 & A sign change through an intermediary costs one delay, about a millisecond; inhibition arriving right after excitation narrows the coincidence window & Rat hippocampal pyramidal neurons in slices: inhibition through one intermediary follows direct excitation after a short delay, and the window in which two excitatory inputs sum to threshold is less than $2\ms$. On the dendrites, where this inhibition is weaker, the window is wider & \cite{Pouille2001} & measured \\
4 & The veto $a\wedge\bar b$~\eqref{eq:veto} with OR and a constant one is functionally complete (Proposition~\ref{prop:gabacomplete}) & Proof in the chapter. A classical result: a network of threshold elements with inhibitory inputs implements any logical expression. A claim about the model; no experiment & \cite{McCullochPitts1943} & theory \\
5 & A delayed veto gives a direction-of-motion detector with optimal speed $v^*=\Delta x/d$~\eqref{eq:vstar} & Rabbit retina: direction selectivity is created by inhibition that cancels excitation for motion in the ``null'' direction. Separate recording of the excitatory and inhibitory inputs of selective cells showed that starburst amacrine cells (\nt{GABA}) inhibit them asymmetrically, only with processes pointing to the ``null'' side. The formula for $v^*$ is the chapter's derivation & \cite{Barlow1965,Fried2002} & measured \\
6 & Shunting inhibition divides the voltage rather than subtracting~\eqref{eq:shunt} (Fig.~\ref{fig:shunt}) & Ohm's law for a divider of three conductances with the chloride equilibrium potential near rest; for a neuron that does not fire & derivation in the chapter & theory \\
7 & On the spike rate the shunt acts by subtraction, and division comes from inhibition together with balanced background noise & Model: in a firing neuron the current through the shunt is nearly constant, and the effect on rate is subtractive. Experiment: background excitatory and inhibitory conductances were injected into rat somatosensory cortex pyramidal neurons (slices) with a dynamic clamp; a simultaneous balanced increase of both divides the slope of the response with no noticeable shift in rate. Review: division by total activity occurs in many parts of the brain & \cite{HoltKoch1997,Chance2002,Carandini2012} & measured \\
8 & For $\beta>1$ mutual inhibition selects a single winner (Proposition~\ref{prop:wta}, \eqref{eq:wta}) & The eigenvalues $-(1\pm\beta)/\tau$ are derived in the chapter. The rates $0.73$ and $0.53$ at $\beta=0.5$ are the fixed point $r_{1,2}=(I_{1,2}-\beta I_{2,1})/(1-\beta^2)$; there is no script in \texttt{models/}. No direct experiment is cited in the chapter & derivation in the chapter & theory \\
9 & Memory reset by a signal through \nt{GABA}; two groups with mutual inhibition form a latch & Follows from Fig.~\ref{fig:bistable} and Proposition~\ref{prop:wta}. No experiment & derivation in the chapter & theory \\
10 & The equilibrium of the \nt{GLUT}--\nt{GABA} loop is stable if the inhibition is strong enough and fast enough (Proposition~\ref{prop:ei}) & Two-population model~\eqref{eq:ei}; conditions (i) and (ii) are the determinant and trace of its matrix, derived in the chapter & \cite{WilsonCowan1972} & theory \\
11 & For $w_{EE}=1.5$ and $w_{EI}w_{IE}=2$, fast inhibition ($\tau_I=2\ms$) holds $E^*=0.67h$, and slow inhibition ($30\ms$) drives oscillations (Fig.~\ref{fig:ei}) & Numerical solution of~\eqref{eq:ei}, script \texttt{figures/make\_ei.py} & calculation & model \\
12 & The cortex works in an inhibition-stabilized regime; the signature: push the \nt{GABA} neurons, and their rate drops~\eqref{eq:isn} & Theory predicted the paradoxical response. Experiment: intracellular recordings in cat primary visual cortex show that when the response is suppressed by a surrounding stimulus, the inhibitory conductance does not rise but falls together with the excitatory one. Optogenetic excitation of PV neurons in mouse visual, somatosensory, and motor cortex lowers the rate of inhibitory neurons, including without a stimulus and under light anesthesia. The coefficient $-1/3$ for the numbers of Fig.~\ref{fig:ei} is the chapter's derivation & \cite{Tsodyks1997isn,Ozeki2009,Sanzeni2020} & measured \\
13 & The loop ``\nt{GLUT} excites \nt{GABA}, \nt{GABA} inhibits \nt{GLUT}'' generates a fast rhythm with a period of $12$--$50\ms$ & Optogenetic stimulation of fast \nt{GABA} neurons in mouse somatosensory cortex at $8$--$200$\,Hz selectively enhances the gamma rhythm ($20$--$80$\,Hz, period $12$--$50\ms$); stimulation of \nt{GLUT} neurons enhances only slower rhythms & \cite{Cardin2009,Buzsaki2012} & measured \\
\bottomrule
\end{longtable}}

\section{Chapter~\ref{sec:code}. Morse Code}

The chapter's limiting estimates are information theory and arithmetic; what has been measured is where noise arises in the channel, how fast the brain works, and what a spike costs, while the question of which code the cortex actually uses remains open.

{\footnotesize
\setlength{\tabcolsep}{3pt}\renewcommand{\arraystretch}{1.15}
\begin{longtable}{>{\raggedright}p{0.5cm}>{\raggedright}p{4.0cm}>{\raggedright}p{5.6cm}>{\raggedright}p{2.3cm}>{\raggedright\arraybackslash}p{1.7cm}}
\toprule
No. & Claim & Experiment and what was measured & Source & Status \\
\midrule\endfirsthead
\toprule
No. & Claim & Experiment and what was measured & Source & Status \\
\midrule\endhead
1 & Rates of cortical \nt{GLUT} neurons range from fractions to tens of spikes per second, and most of the time neurons work near the lower end & Cell-attached recording (neuron selection independent of its activity) in the auditory cortex of awake rats: at any moment fewer than $5\,\%$ of neurons fire above $20$ spikes per second; some neurons have a background rate below $0.01$ per second; the rate distribution is lognormal & \cite{Hromadka2008} & measured \\
2 & The capacity of a spike channel is $R_{\max}\approx r\log_2\frac{e}{r\Delta t}$~\eqref{eq:spikeentropy}, and almost all of it lies in spike timing (Proposition~\ref{prop:capacity}) & Entropy of a binary string of bins of length $\Delta t$. The chapter's numbers: $8$ bits per spike at $r=10$ per second and $\Delta t=1\ms$, about $5$ bits at $\Delta t=10\ms$, fewer than $2$ bits per second for a spike counter & \cite{MacKay1952} & theory \\
3 & Sensory neurons transmit one to several bits per spike, in the best cases up to half the limit & Stimulus reconstruction from the spikes of sensory neurons whose stimulus is known; a summary of measurements in the book & \cite{Rieke1997} & measured \\
4 & The spike generator is precise; precise timing is set by fast transients of the input & Rat cortical neurons in slices: for a repeated current with fast fluctuations, spikes repeat with a precision better than $1\ms$; for a constant current the times drift & \cite{Mainen1995} & measured \\
5 & The synapse is unreliable: more than half of the spikes do not reach the recipient because of random release & Minimal stimulation of inputs to hippocampal pyramidal neurons (slices): more than half of the spikes evoke no response. Fluctuations of the stimulation threshold, conduction failures at axon branch points, and low experimental temperature were ruled out; what remains is probabilistic release & \cite{AllenStevens1994} & measured \\
6 & The spike train in living cortex looks random: intervals are scattered as in a Poisson stream, and the variance of the spike count is close to the mean; part of the ``noise'' is a message about the animal's movements & Recordings in visual areas V1 and MT of awake monkeys: coefficient of variation of the intervals about $1$, ratio of spike-count variance to mean as for a random process; a model summing independent small inputs gives much less variability. In mouse visual cortex a considerable part of the variability is predicted from video of the animal's own movements & \cite{Softky1993,Stringer2019} & measured \\
7 & The rate code is slow: error $1/\sqrt{rT}$~\eqref{eq:rateerror}, while the brain recognizes a picture in $\sim150\ms$ & The formula is Poisson statistics, the chapter's derivation. Experiment: people decided whether an unfamiliar photograph flashed for $20\ms$ contained an animal; the brain potentials for ``yes'' and ``no'' diverge after about $150\ms$. Splitting this time into ten stages of $10$--$20\ms$ is the chapter's estimate, not a measurement & \cite{Thorpe1996} & measured \\
8 & Averaging over a group is limited by correlation: the error falls by no more than a factor of $1/\sqrt c$~\eqref{eq:correlated} (Proposition~\ref{prop:pooling}) & Pairs of monkey MT neurons recorded simultaneously: the spike-count correlation coefficient averages $0.12$, and theoretical analysis shows that even such a correlation limits the gain from a group. Theory: only the part of the correlation that looks like the signal sets the limit. A model of the opposing position: the rate of a group of $50$--$100$ neurons is read within one interspike interval, $10$--$50\ms$, and the timing of individual spikes is not used in the cortex & \cite{Zohary1994,MorenoBote2014,ShadlenNewsome1998} & measured \\
9 & A number can be written in the first-spike time~\eqref{eq:latency}, in the spike order of a group, or in the phase of a local rhythm & Salamander retina: the spatial structure of a briefly flashed image is recorded in the relative latencies of the first spikes of the output neurons, almost independently of contrast. Rat hippocampal place cells: while passing through ``their'' place, spikes shift earlier and earlier in the theta cycle ($7$--$12$\,Hz), a shift of $100^\circ$ to $355^\circ$. How much the cortex uses spike timing is an open question (row~8) & \cite{Gollisch2008,OKeefe1993} & measured \\
10 & Which code gets read is decided by the recipient's time constant~\eqref{eq:reader} (Proposition~\ref{prop:reader}); in active cortex neurons are closer to coincidence detectors & Formula~\eqref{eq:reader} is the chapter's derivation. Review of in vivo recordings: in active cortex the membrane conductance is several times higher than at rest, and the time constant correspondingly shorter. That the cortex switches codes in precisely this way has not been shown directly & \cite{Destexhe2003} & hypothesis \\
11 & A depressing synapse transmits the change in rate, essentially $\Delta r/r$~\eqref{eq:depression} (Fig.~\ref{fig:depression}) & Paired recordings of cortical pyramidal neurons: the rate of depression is set by the release probability; with slow depression the response reflects rate, with fast depression coincidences. A model based on measured depression in rat cortex: equal relative rate changes on fast and slow inputs give equal responses, that is, the synapse transmits $\Delta r/r$. The numbers $1.7\to2.2$, $6.7$, and $90\ms$ are the chapter's calculation & \cite{Tsodyks1997,Abbott1997depression} & measured \\
12 & A burst is a ``dash'': its probability of being lost, $(1-p)^k$~\eqref{eq:burst}, is small, and facilitation lowers it further & Review: at unreliable synapses single spikes are often lost, while bursts are transmitted reliably thanks to facilitation; a burst induces long-term synaptic changes. That the brain reads a burst as a separate symbol is a hypothesis & \cite{Lisman1997,AllenStevens1994} & hypothesis \\
13 & Fast \nt{GABA} neurons set the rhythm and ``punctuation''; one more class inhibits the inhibitors and opens the gates (Proposition~\ref{prop:punctuation}) & Review of the properties of classes of cortical \nt{GABA} neurons. Optogenetics: rhythmic excitation of fast \nt{GABA} neurons evokes the gamma rhythm, and the response to a stimulus depends on the phase of the cycle. In mouse visual cortex PV neurons inhibit one another, SST neurons inhibit all the other classes, and VIP neurons inhibit mainly SST. Excitation of VIP neurons in awake mice releases \nt{GLUT} neurons from inhibition; they are switched on by a reinforcement signal. The division of roles as a general rule is a hypothesis & \cite{Tremblay2016,Cardin2009,Pfeffer2013,Pi2013} & hypothesis \\
14 & Energy limits the mean rate to about one spike per second~\eqref{eq:energy} & Energy budget of rodent gray matter from anatomical and physiological data: spikes and synaptic currents take $47\,\%$ and $34\,\%$ of the signaling energy; no more than $\sim15\,\%$ of neurons can be active at once. For the human cortex: no more than $\sim1\,\%$ of neurons can be strongly active at once. The figures of $\sim10^9$ ATP per spike and $\sim10$\,W for signaling are the chapter's estimate based on these calculations & \cite{Attwell2001,Lennie2003} & theory \\
15 & The code is sparse and tuned for the greatest number of bits per joule; the cortex trades precision for energy (Proposition~\ref{prop:economy}) & The efficient-coding hypothesis. Support: in food-restricted mice, AMPA receptor conductance and synaptic ATP use in the visual cortex fall ($-29\,\%$), the firing rate is preserved, and orientation tuning broadens by $32\,\%$ & \cite{Barlow1961,Padamsey2022} & hypothesis \\
\bottomrule
\end{longtable}}

\section{Chapter~\ref{sec:serocomputer}. \nt{SERO} neurons}

The properties of \nt{SERO} neurons themselves (rhythm, autoinhibition, sign set by the receptor, subsystems) are measured directly; the chapter's computations (regulator, filter, logic) follow from its equations; $\tau_c$, the diffusion coefficient of serotonin and the number of release sites are estimates; the loop's role as a level regulator in the brain is a hypothesis (in the raphe nucleus autoinhibition is point to point and causes only brief pauses), and the role of \nt{SERO} as the horizon is a hypothesis with behavioral experiments in its favor.

{\footnotesize
\setlength{\tabcolsep}{3pt}\renewcommand{\arraystretch}{1.15}
\begin{longtable}{>{\raggedright}p{0.5cm}>{\raggedright}p{4.0cm}>{\raggedright}p{5.6cm}>{\raggedright}p{2.3cm}>{\raggedright\arraybackslash}p{1.7cm}}
\toprule
No. & Claim & Experiment and what was measured & Source & Status \\
\midrule\endfirsthead
\toprule
No. & Claim & Experiment and what was measured & Source & Status \\
\midrule\endhead
1 & The sign of serotonin's action is chosen by the recipient's receptor (Proposition~\ref{prop:receptor}) & Serotonin receptors are divided into families by pharmacology and mechanism: 5-HT$_{1A}$ opens potassium channels through a G protein and inhibits the cell; 5-HT$_{2A}$ and the ionotropic 5-HT$_3$ excite. One transmitter gives opposite responses in different cells. & \cite{BarnesSharp1999} & measured \\
2 & In the waking state a \nt{SERO} neuron steadily fires a few spikes per second & Recordings of dorsal raphe neurons in freely moving cats: slow rhythmic firing of $2.8\pm0.2$ spikes/s in quiet waking; in slow-wave sleep the rate falls by $34$--$68\,\%$, in REM sleep by $98\,\%$. In anesthetized rats, $1.7\pm0.5$ spikes/s, very regular. & \cite{TrulsonJacobs1979, GagnonParent2014, JacobsAzmitia1992} & measured \\
3 & A \nt{SERO} neuron is a pacemaker: it needs no push for each spike, but it needs a steady background drive (in a slice, noradrenaline through $\alpha_1$ receptors; in the organism, from \nt{NORA}) & In a brain slice the cells fire slowly and steadily, with an afterhyperpolarization of $200$--$800\ms$ after each spike. Fewer cells fire spontaneously in a slice than in the organism: the steady rhythm needs a tonic noradrenaline input through $\alpha_1$ receptors, and blocking them abolishes it. & \cite{VandermaelenAghajanian1983} & measured \\
4 & Almost all receptors are metabotropic; the response is slow, rising and decaying within about a second & All receptor families except 5-HT$_3$ are coupled to G proteins. In a slice, evoked serotonin release produces an inhibitory current through 5-HT$_{1A}$ that rises and decays within a second. & \cite{BarnesSharp1999, CourtneyFord2016} & measured \\
5 & In target regions serotonin is released into the volume, not only into the cleft; in the raphe nucleus itself, inhibition of neighbors is point to point & Fast-scan cyclic voltammetry in slices: release per spike is the same in a region with synapses and in the raphe nucleus, where there are almost no synapses; it does not depend on blocking uptake or receptors; there are far fewer molecules per terminal than transporters and receptors, and the extracellular concentration matches the affinity of the main receptor. In the raphe nucleus, inhibition through 5-HT$_{1A}$ is point to point: the transporter keeps serotonin from accumulating outside the cleft. & \cite{Bunin1998, CourtneyFord2016} & measured \\
6 & The concentration in~\eqref{eq:seroneuron} decays because of reuptake; $\tau_c$ of order a second is an estimate & Voltammetry in the living brain: extracellular serotonin is limited primarily by reuptake and degradation, not by synthesis. The cited papers contain no direct measurement of $\tau_c$; its order of magnitude is the chapter's estimate. & \cite{Hashemi2012serotonin} & measured; $\tau_c$ is an estimate \\
7 & NOT and NOR with a single pacemaker; completeness of \nt{SERO} logic (Proposition~\ref{prop:serocomplete}, Fig.~\ref{fig:serogates}) & Follows from~\eqref{eq:seroneuron}: an inhibitable pacemaker is a NOR gate, and NOR is functionally complete. There is no experiment in which logic has been built from \nt{SERO} neurons; the condition of separate volumes does not hold in the brain. & derived in the chapter & theory \\
8 & Serotonin inhibits the \nt{SERO} neuron itself through receptors on it & In anesthetized rats, 5-HT$_{1A}$ agonists given intravenously and iontophoretically inhibit the firing of dorsal raphe neurons with the same dose--response relation as serotonin itself; 5-HT$_{1B}$ agonists have almost no effect. In a slice, endogenous serotonin from neighboring cells produces a current through 5-HT$_{1A}$ and a brief pause in firing without changing the mean rate. & \cite{Sprouse1987, CourtneyFord2016} & measured \\
9 & In the model the loop through a common volume is a level regulator: variation and the time constant shrink by a factor of $1+G$~\eqref{eq:feedback}; that the loop works this way in the brain is a hypothesis & The classical formula of an amplifier with negative feedback, derived again in the chapter. The loop gain $G$ of the \nt{SERO} system has not been measured. Measured: in a slice, autoinhibition is point to point, causes a brief pause and does not change the mean rate. & \cite{Black1934feedback, CourtneyFord2016} & theory; in the brain, hypothesis \\
10 & The concentration is a moving average of the rate, a low-pass filter~\eqref{eq:lowpass}, Fig.~\ref{fig:lowpass} & Solution of the linear equation~\eqref{eq:seroneuron}; the figure is computed from this formula with $\tau_c=1\,\text{s}$. There is no separate model in \texttt{models/}. & derived in the chapter & theory \\
11 & The tortuosity of the gaps slows the diffusion of a small molecule by a factor of about $2.6$ & Point-source method: iontophoresis of tetramethylammonium and recording of its concentration with an ion-selective electrode in slices and in the living brain. Tortuosity $\lambda\approx1.6$, so the effective $D$ is $\lambda^2\approx2.6$ times smaller than the free one; the extracellular volume fraction is about $0.2$. The diffusion coefficient of serotonin itself in tissue was not measured in these studies; in the chapter it is an estimate. & \cite{Sykova2008} & measured \\
12 & The length of an independent ``wire'' is $\ell\sim\sqrt{D\tau}\approx20$~\textmu m~\eqref{eq:diffusionlength}; the number of wires is limited by the number of such regions (Proposition~\ref{prop:volumewires}) & The diffusion law with the estimates of $D$ and $\tau$ substituted (rows~6 and~11); Proposition~\ref{prop:volumewires} is a consequence. There is no direct experiment counting the independent channels of volume transmission. & derived in the chapter & theory \\
13 & The output of a single \nt{SERO} neuron is spread over huge parts of the brain; the number of release sites is estimated at $\sim10^5$ & Complete reconstruction of single dorsal raphe neurons in the rat: all ascending axons branch heavily, the total axon length of one cell reaches $18.7$~cm, and the branches of one cell reach the striatum, prefrontal cortex and amygdala. The number of release sites per cell has not been counted directly. & \cite{GagnonParent2014} & measured; $10^5$ is an estimate \\
14 & \nt{SERO} neurons fall into several subsystems with different outputs and responses & Viral-genetic labeling in mice: neurons projecting to the amygdala and to the frontal cortex hardly overlap in their branching, receive different inputs and respond oppositely to an aversive stimulus; activating and silencing each group changes behavior in different ways. & \cite{Ren2018} & measured \\
15 & Patience condition $D<\tau_\gamma\ln(R/r_0)$~\eqref{eq:patience} with exponential discounting; with the measured, hyperbolic discounting, $D<\tau_\gamma(R/r_0-1)$ & Follows from discounting by $e^{-D/\tau_\gamma}$ or $1/(1+D/\tau_\gamma)$; derived in the chapter. Monkeys choosing at delays of seconds: discounting is closer to a hyperbola than to an exponential; the response of \nt{DOPA} neurons to the cue of a delayed reward also falls hyperbolically with delay. & derived in the chapter; \cite{KobayashiSchultz2008} & theory \\
16 & \nt{SERO} sets the horizon $\tau_\gamma$ (Section~\ref{sec:horizon}) & Optogenetic activation of dorsal raphe \nt{SERO} neurons in mice lengthens waiting for a delayed reward and increases persistence in foraging for a reward that has become scarce. In humans, lowering serotonin (a tryptophan-free diet) makes the discounting of delayed reward steeper. How the concentration becomes $\tau_\gamma$ is unknown. & \cite{Doya2002, Miyazaki2014, Lottem2018, Schweighofer2008} & hypothesis \\
\bottomrule
\end{longtable}}

\section{Chapter~\ref{sec:dofa}. Learning with \nt{DOPA}}

The mechanisms of the synapse (packets, coincidence detector, calcium, plasticity windows) and the behavior of dopamine as a prediction error are measured directly; that the rules lead to the gradient, and what broadcasting costs, are theorems; the outcome of learning to choose is computed with the book's model; the circuit that computes the error is a hypothesis.

{\footnotesize
\setlength{\tabcolsep}{3pt}\renewcommand{\arraystretch}{1.15}
\begin{longtable}{>{\raggedright}p{0.5cm}>{\raggedright}p{4.0cm}>{\raggedright}p{5.6cm}>{\raggedright}p{2.3cm}>{\raggedright\arraybackslash}p{1.7cm}}
\toprule
No. & Claim & Experiment and what was measured & Source & Status \\
\midrule\endfirsthead
\toprule
No. & Claim & Experiment and what was measured & Source & Status \\
\midrule\endhead
1 & Backpropagation of error in the form used in artificial networks has not been found in the brain & There is no direct experiment: this is a claim of absence. Reviews analyze what the algorithm requires (feedback connections mirroring the forward ones, a separate phase) and which biological approximations to it have been proposed. & \cite{Lillicrap2020, WhittingtonBogacz2019} & hypothesis \\
2 & The weight factors into the number of release sites, the release probability and the response to one packet: $w=N_s\,p\,A_1$~\eqref{eq:quantal} & Frog neuromuscular junction with reduced release: the recipient's response fluctuates in steps that are multiples of the spontaneous miniature response to one packet; the number of packets per spike obeys the Poisson law. & \cite{delCastillo1954} & measured \\
3 & The recipient's receptor is a coincidence detector; calcium triggers the weight change & Patch clamp of mouse neurons: Mg$^{2+}$ ions block the glutamate-opened channel at the resting potential and leave it upon depolarization. Hippocampal slices: a calcium chelator in the recipient blocks long-term potentiation, while light-triggered release of calcium inside the cell by itself strengthens transmission. & \cite{Nowak1984, Malenka1988calcium} & measured \\
4 & Either side carries out the decision: $A_1$ grows in the recipient, $p$ changes in the sender & Glutamate uncaged by light over a single spine causes rapid and lasting growth of the spine and of the current through its receptors; potentiation is accompanied by insertion of AMPA receptors. Depolarizing the recipient temporarily suppresses release by the sender; the effect is carried by endocannabinoids traveling back across the cleft (shown at \nt{GABA} inputs). Short-term depression depends on the sender's release probability. & \cite{Matsuzaki2004, Malinow2002, WilsonNicoll2001, Tsodyks1997} & measured \\
5 & A synapse strengthens when the sender and the recipient are active together~\eqref{eq:hebb} & Anesthetized rabbit: after brief trains of spikes along the hippocampal input pathway, the response is enhanced for $30$~min to $10$~h. In a hippocampal slice, the recipient's spikes strengthen only the synapses that are active at the same moment. & \cite{Bliss1973LTP, Kelso1986hebb} & measured \\
6 & Rule~\eqref{eq:oja} finds the principal eigenvector of the input correlations (Proposition~\ref{prop:oja}) & A theorem; the proof is in the chapter. There is no experiment in which a neuron has learned the principal component of its inputs; the mechanism that bounds weight growth in the synapse has not been established. & \cite{Oja1982} & theory \\
7 & Timing-based Hebb: a window of about $\pm20\ms$ (Fig.~\ref{fig:stdp}); chains grow from repeated sequences & Pairs of hippocampal neurons in culture: a recipient spike within $20\ms$ after a sender spike gives lasting strengthening, within $20\ms$ before, weakening; both depend on NMDA receptors. The ratio $A_-=0.6A_+$ in the figure is illustrative. That chains grow this way in a network has not been measured directly. & \cite{BiPoo1998} & measured \\
8 & Eligibility trace~\eqref{eq:threefactor}: dopamine arriving $1$--$2\,\text{s}$ after joint activity strengthens the connection, later dopamine does not (Proposition~\ref{prop:threefactor}) & Striatal slices, separate optogenetic stimulation of glutamate and dopamine inputs: the spine grows only if dopamine arrives $0.3$--$2\,\text{s}$ after the glutamate input; the window is set by the fast rise and decay of cAMP. In the cortex, pairing leaves separate traces for strengthening and weakening, which monoamine receptors convert into long-term changes within a window of order seconds. & \cite{Yagishita2014, He2015eligibility} & measured \\
9 & Windows lasting seconds also occur without dopamine: the third signal is strong excitation of the recipient & Awake mouse, area CA1: a single plateau event in the recipient strengthens inputs that arrived seconds before and after it, and creates a place field in a single pass. & \cite{Bittner2017} & measured \\
10 & The mean step of the three-factor rule follows the reward gradient~\eqref{eq:perturbation} (Proposition~\ref{prop:gradient}) & The gradient theorem for learning from random perturbations; derived in the chapter for small noise. & \cite{Williams1992} & theory \\
11 & Noise is search: the network learns from its own random deviations & Juvenile zebra finches: inactivating the output nucleus of the basal ganglia circuit abruptly and reversibly removes song variability. Adult zebra finches: noise is played on renditions of a syllable with pitch on one side of the mean, and the birds quickly shift the pitch the other way, learning from their own variability. & \cite{Olveczky2005, TumerBrainard2007} & measured \\
12 & A two-way choice is learned with $\tau_e=1\,\text{s}$ and $\delta=R-\bar R$; it is not learned with $\tau_e=50\ms$; without subtraction the weights grow without bound, and at a high learning rate the network can lock in the worse choice & \texttt{models/dofa\_learning.py}, $\eta=0.05$, $500$ trials, $6$ seeds. $\tau_e=1\,\text{s}$: fraction of $A$ choices $0.65$--$0.79\to0.96$--$1.00$, weights $0.85$--$1.33$. $\tau_e=50\ms$: $0.38$--$0.55$, weights unchanged. $\delta=R$: winner's weight $860$--$1190$. $\delta=R$, $\eta=0.5$, $3$ seeds: one locked onto $B$ ($0.04\to0$). The script prints all four cases; the rerun matches the chapter. & \texttt{models/\allowbreak dofa\_\allowbreak learning.py} & model \\
13 & The learning time from one shared signal grows in proportion to the number of noisy neurons (Proposition~\ref{prop:broadcastcost}) & Learning curves of a linear network: node perturbation is slower than the exact gradient by a factor equal to the number of output neurons, and weight perturbation is slower by a further factor equal to the number of inputs. & \cite{Werfel2005} & theory \\
14 & \nt{DOPA} outputs go primarily to action-selection regions & Complete axon reconstruction of single nigrostriatal \nt{DOPA} neurons in the rat: the branches of one cell cover $0.45$--$5.7\,\%$ of the striatal volume ($2.7\,\%$ on average), with almost no branches outside the striatum. & \cite{Matsuda2009} & measured \\
15 & The cortex learns to predict its input, and the error is computed on the spot & Review: sensory cortex shows responses to mismatch between expected and actual input. That this is a general learning rule of the cortex has not been proven. & \cite{Keller2018} & hypothesis \\
16 & Dopamine behaves like the error~\eqref{eq:ctd}: burst, transfer to the cue, dip (Fig.~\ref{fig:ctdcases}) & Monkey \nt{DOPA} neurons: a burst to unexpected juice; after learning, a burst to the cue and no response to the promised juice; a dip when the juice is omitted. The continuous form~\eqref{eq:ctd} is theory. & \cite{Schultz1997, Doya2000} & measured \\
17 & A circuit that computes $\dd V/\dd t$ and subtracts the expectation & Mice, recording and optogenetics in the ventral tegmental area: \nt{DOPA} neurons subtract the expectation from the reward; neighboring \nt{GABA} neurons inhibit them when reward is expected, and exciting or inhibiting these \nt{GABA} neurons changes the error. How the derivative is computed has not been shown. & \cite{Eshel2015arithmetic} & hypothesis \\
18 & The \nt{DOPA} neuron is a pacemaker; dips are shallower than bursts & In a slice, \nt{DOPA} neurons fire spontaneously and very regularly, riding on a slow pacemaker potential. In the monkey, the rate follows a positive error quantitatively but conveys a negative error only weakly. & \cite{GraceOnn1989, Bayer2005} & measured \\
19 & Actor and critic (Fig.~\ref{fig:actorcritic}) & The algorithm comes from reinforcement-learning theory. In humans (fMRI), the ventral striatum behaves like a critic both with and without choice, while the dorsal striatum does so only when actions must be chosen. & \cite{SuttonBarto2018, ODoherty2004actor} & hypothesis \\
20 & The sign of $\delta$ in the rule is flipped in neurons with the second receptor family & Striatal slices: in neurons with D1 receptors, pairing produces strengthening that requires D1; in neurons with D2, weakening that requires D2. & \cite{Shen2008} & measured \\
\bottomrule
\end{longtable}}

\section{Chapter~\ref{sec:dofaalt}. Other theories of \nt{DOPA}}

Each of the extended pictures rests on its own direct measurements of dopamine (the spread of reversal points, responses to unpleasant and novel events, the slow ramp, responses to a flavor switch), but the formulas that explain them are theories, and between two of them the experiments so far contradict one another.

{\footnotesize
\setlength{\tabcolsep}{3pt}\renewcommand{\arraystretch}{1.15}
\begin{longtable}{>{\raggedright}p{0.5cm}>{\raggedright}p{4.0cm}>{\raggedright}p{5.6cm}>{\raggedright}p{2.3cm}>{\raggedright\arraybackslash}p{1.7cm}}
\toprule
No. & Claim & Experiment and what was measured & Source & Status \\
\midrule\endfirsthead
\toprule
No. & Claim & Experiment and what was measured & Source & Status \\
\midrule\endhead
1 & The asymmetric rule~\eqref{eq:asymmetric} converges to the point~\eqref{eq:expectile}; for $\kappa=1/2$ it is the mean, and for a drop with probability $1/2$, $V=\kappa$ (Fig.~\ref{fig:expectiles}) & The point~\eqref{eq:expectile} is an expectile, the solution of a least-squares problem with different weights for positive and negative deviations; for the chapter's example the derivation is in the chapter itself. & \cite{NeweyPowell1987}; derived in the chapter & theory \\
2 & Different \nt{DOPA} neurons have different reversal points, ordered by asymmetry (Proposition~\ref{prop:expectile}) & Mice, optogenetically tagged \nt{DOPA} neurons of the ventral tegmental area, rewards of different sizes: the point where a burst turns into a dip differs between neurons; it is higher in neurons with greater response asymmetry; the shape of the reward distribution can be reconstructed from the group's responses. That this is the main function of the diversity is a hypothesis. & \cite{Dabney2020} & measured \\
3 & Some \nt{DOPA} neurons respond to unpleasant events with a burst & Monkeys: some \nt{DOPA} neurons are excited by a reward cue and inhibited by a cue for an air puff to the eye, others are excited by both; the groups lie in different parts of the midbrain. & \cite{Matsumoto2009, BrombergMartin2010} & measured \\
4 & The magnitude of the error, or novelty, sets the learning rate & The cited papers contain no direct experiment in which a salience signal in \nt{DOPA} neurons changes the learning rate; the review proposes the role of an ``attention, this matters'' signal. & \cite{BrombergMartin2010} & hypothesis \\
5 & A separate wire for the danger axis & Mice: \nt{DOPA} neurons projecting to the tail of the striatum respond to novel and intense stimuli, not to reward; ablating them weakens avoidance of a novel object. & \cite{Menegas2018} & measured \\
6 & The optimal action time $\tau_a^*=\sqrt{C/\bar R}$~\eqref{eq:vigor} & Minimum of the sum $C/\tau_a+\bar R\tau_a$; derived in the chapter. In the original model, the speed of actions is set by the cost of time, equal to the average reward. & \cite{Niv2007}; derived in the chapter & theory \\
7 & The slow dopamine level carries $\bar R$ and sets the speed of actions (Proposition~\ref{prop:multiplex}) & Rats, microdialysis and voltammetry in the nucleus accumbens: the dopamine level tracks the reward rate and the speed of actions from minute to minute. In humans, L-DOPA strengthens the dependence of reaction time on the average reward rate. That the slow level equals precisely $\bar R$ has not been proven. & \cite{Hamid2016, Beierholm2013vigor} & hypothesis \\
8 & The slow level combines two sources: release changes at the outputs without a change in spike rate, but a slow ramp also occurs in the rate itself & Rats, one task: release in the nucleus accumbens follows the changing reward expectation, while the spike rate of identified \nt{DOPA} neurons in the ventral tegmental area does not. Mice in a virtual corridor (row~10): a smooth ramp on approaching reward is present in the spikes of identified \nt{DOPA} neurons as well. & \cite{Mohebi2019, Kim2020} & measured \\
9 & Dopamine rises smoothly while the animal approaches a distant reward & Rats in a maze, voltammetry in the striatum: the dopamine concentration builds up over seconds as the goal approaches; the slope depends on the distance and on the reward size. & \cite{Howe2013ramp, Hamid2016} & measured \\
10 & The ramp is the error $\delta$ as the estimate is refined~\eqref{eq:ramp}; a jump in the corridor gives a burst (Fig.~\ref{fig:ramp}) & The formula and the value $\delta=0.75\,V$ per second are derived in the chapter. Experiment: mice in a virtual corridor, instantaneous teleportation closer to the goal; somatic spikes, somatic and axonal calcium, and the concentration in the striatum respond like an error, not like the expectation $V$. Which of the two explanations holds at all outputs is debated. & \cite{Kim2020} & measured \\
11 & Looking back: dopamine reports a measure of causal link~\eqref{eq:retro} & Mice, dopamine release in the nucleus accumbens in experiments where this theory and the error~\eqref{eq:ctd} predict different things: in several experiments release followed the former. & \cite{Jeong2022} & hypothesis \\
12 & The dopamine response gradually shifts from reward to cue during learning & Mice, signal from \nt{DOPA} neuron outputs in the ventral striatum over the course of learning: the burst gradually moves from the moment of reward to the moment of the cue, as learning by~\eqref{eq:ctd} requires. & \cite{Amo2022} & measured \\
13 & \nt{DOPA} neurons respond to a change in the reward's flavor at the same value & Rats, recordings from the ventral tegmental area: replacing a juice with another of equal value evokes a response, even though the error~\eqref{eq:ctd} is zero. & \cite{Takahashi2017} & measured \\
14 & Artificial bursts teach an association between two neutral stimuli & Rats, sensory preconditioning of two stimuli without reward: optogenetic activation of \nt{DOPA} neurons at the transition from the first stimulus to the second creates the association, and inhibition prevents it from being learned. & \cite{Sharpe2017} & measured \\
15 & A vector of errors per feature~\eqref{eq:vectortd} & A theory of feature-prediction errors; there is no direct test of the vector form. & \cite{Gardner2018} & hypothesis \\
16 & Different \nt{DOPA} neurons in the same task carry different quantities & Mice in a virtual corridor, two-photon imaging of \nt{DOPA} neurons: some are related to reward, others to speed and turning, still others to cues and choice accuracy. & \cite{Engelhard2019} & measured \\
17 & A bundle instead of a wire (Proposition~\ref{prop:bundle}) & Summary of rows~2--16: each decomposition is measured separately; there is no experimental evidence for a single picture in which all of them are parts of one signal. & --- & hypothesis \\
\bottomrule
\end{longtable}}

\section{Chapter~\ref{sec:nora}. \nt{NORA}}

The organization of the \nt{NORA} system, its two modes, its link to the pupil (not only through \nt{NORA}), the rise in signal-to-background ratio and the inverted U in behavior are measured directly; the multiplicative gain $g$ is a model; that gain switches the choice mode and acts as an inverse temperature are conclusions of the chapter; the benefit of adjusting the gain is computed with the book's model; ``\nt{NORA} as the exploration knob'' and ``\nt{NORA} as an interrupt'' are hypotheses.

{\footnotesize
\setlength{\tabcolsep}{3pt}\renewcommand{\arraystretch}{1.15}
\begin{longtable}{>{\raggedright}p{0.5cm}>{\raggedright}p{4.0cm}>{\raggedright}p{5.6cm}>{\raggedright}p{2.3cm}>{\raggedright\arraybackslash}p{1.7cm}}
\toprule
No. & Claim & Experiment and what was measured & Source & Status \\
\midrule\endfirsthead
\toprule
No. & Claim & Experiment and what was measured & Source & Status \\
\midrule\endhead
1 & A human has a few tens of thousands of \nt{NORA} neurons, almost all in one group & Serial sections of the human brainstem stained for tyrosine hydroxylase: $53\,900$ cells in the locus coeruleus and another $6\,260$ in the adjacent subnucleus. & \cite{Baker1989LC} & measured \\
2 & The outputs spread through almost the whole brain, but parts of the group partly serve different regions & Viral-genetic labeling in mice: \nt{NORA} neurons of the locus coeruleus receive inputs from wide areas of the brain and send outputs to wide areas of the brain and spinal cord. In rats: a group projecting to the amygdala enhances fear learning, and a separate group projecting to the prefrontal cortex enhances its extinction. & \cite{Schwarz2015LC, Uematsu2017LC, Poe2020} & measured \\
3 & Tonic mode: steady spikes at rates from a fraction of a spike to a few per second, changing with wakefulness & Rats, recordings of locus coeruleus neurons across the sleep--wake cycle: the rate is highest in waking, lower in slow-wave sleep, and almost zero in REM sleep; rate changes precede state transitions. & \cite{AstonJonesBloom1981} & measured \\
4 & Phasic mode: a short burst to a task-relevant event, at roughly the moment of decision & Monkeys, a task of detecting a rare target: all locus coeruleus neurons respond with a burst only to the target, $\sim90\ms$ after it and $\sim200\ms$ before the hand response; bursts are weaker when performance is poor. In a two-choice task the burst is more tightly locked to the response than to the stimulus, and it occurs on both correct and error trials. & \cite{AstonJones1994vigilance, Clayton2004LC} & measured \\
5 & Pupil diameter is an imprecise test point for \nt{NORA}: it follows \nt{NORA}, \nt{ACET} and other regions & Monkeys: locus coeruleus spikes, down to single spikes of one cell, predict pupil dilation; a similar but later relation exists in the colliculi and the cingulate cortex. Mice: pupil fluctuations follow the activity of both noradrenergic and cholinergic outputs in the cortex. & \cite{Joshi2016, Reimer2016} & measured \\
6 & Noradrenaline suppresses background activity more than evoked activity; the multiplicative gain~\eqref{eq:gain} is a model that captures this effect & Awake monkeys, auditory cortex, iontophoresis of noradrenaline: background activity of neurons is suppressed more than the response to conspecific calls, so the signal-to-noise ratio rises. The multiplicative sigmoid~\eqref{eq:gain} is taken from a network model; a literal multiplication of the slope has not been measured. & \cite{Foote1975NE, ServanSchreiber1990} & measured; $g$ is a model \\
7 & Gain raises $d'$ only against noise after the amplifier~\eqref{eq:gainsnr} (Proposition~\ref{prop:gainnoise}) & Derived in the chapter; in the network model, gain improves signal detection. There is no experiment that separates noise before and after the amplifier and tests this distinction. & derived in the chapter; \cite{ServanSchreiber1990} & theory \\
8 & The boundary $g\beta=1$ switches the choice network from ``deliberating'' to ``decided''~\eqref{eq:wtagain} (Proposition~\ref{prop:gainwta}, Fig.~\ref{fig:pitchfork}) & Linear analysis of two mutually inhibiting groups; derived in the chapter. There are no direct measurements of this threshold in the brain. & derived in the chapter & theory \\
9 & The gain is the inverse temperature of choice~\eqref{eq:softmax} & With logistic noise after the amplifier, the choice probability is softmax with $T=\sigma/g$; derived in the chapter. & derived in the chapter & theory \\
10 & \nt{NORA} sets how much to explore: high tonic activity means small effective $g$ (exploration), a phasic burst at the moment of decision means large $g$ (decision) & Humans: the baseline pupil is wider before a random choice than before choosing the best known option; people with wider pupils are more inclined to explore. Rats: switching into a random-choice mode is controlled by \nt{NORA} input to the anterior cingulate cortex. That the burst raises the effective $g$ has not been measured directly. & \cite{JepmaNieuwenhuis2011, Tervo2014stochastic, AstonJones2005} & hypothesis \\
11 & Reward depends on $g$ as an inverted U; the best $g$ is smaller in a fast-changing world (Proposition~\ref{prop:explore}, Fig.~\ref{fig:invertedU}) & \texttt{models/nora\_explore.py}, $60$ runs of $4000$ trials. A change once per $1000$ trials: best $g=16$, fraction $0.976$; once per $20$: $g=8$, fraction $0.886$; at $g=0.5$, $0.77$. The rerun matches the chapter. & \texttt{models/\allowbreak nora\_\allowbreak explore.py} & model \\
12 & The inverted U in behavior: best performance at moderate tonic activity with sharp bursts & Monkeys, the same task as in row~4: during periods of good performance the tonic rate is moderate and bursts to the target are sharp; at a high tonic rate there are no bursts and more false responses. Changes in tonic and evoked activity follow fluctuations in performance. & \cite{Usher1999LC, AstonJones2005} & measured \\
13 & \nt{NORA} reports unexpected uncertainty, \nt{ACET} expected uncertainty & A theory of Bayesian inference with two kinds of uncertainty; the cited papers contain no direct experiment separating these signals by transmitter. & \cite{YuDayan2005} & hypothesis \\
14 & After an abrupt change people learn faster, and the pupil dilates & Humans, a prediction task with sudden changes: brief pupil changes track the estimated probability of a change and, together with the baseline pupil, predict how strongly new data change the estimate (the learning rate); pupil changes unrelated to light and task alter this rate. That the pupil here reflects \nt{NORA} specifically is an inference from the indirect data of row~5. & \cite{Nassar2012} & measured \\
15 & The phasic burst works as an interrupt: the network resets its state & There is no direct experiment in which a \nt{NORA} burst resets the network's state; only bursts to important events have been measured (row~4). & \cite{BouretSara2005} & hypothesis \\
16 & Surprise-driven adjustment of $g$ or of the learning rate is hardly better than the best constant settings & \texttt{models/nora\_explore.py}, a world with epochs of $1000$ trials (a change once per $1000$ and once per $20$). Best constant $g=8$: $0.925$; adjusting $g$: up to $0.927$; adjusting the learning rate: up to $0.926$; constant learning rate $0.4$: $0.928$. The rerun matches the chapter. & \texttt{models/\allowbreak nora\_\allowbreak explore.py} & model \\
\bottomrule
\end{longtable}}

\section{Chapter~\ref{sec:acet}. Adding \nt{ACET}}

The three actions of acetylcholine are measured in slices and in the living brain, the conflict between writing and reading is shown with the book's model, and the idea that \nt{ACET} serves as a write-enable line and reports expected uncertainty is a hypothesis with partial experimental support.

{\footnotesize
\setlength{\tabcolsep}{3pt}\renewcommand{\arraystretch}{1.15}
\begin{longtable}{>{\raggedright}p{0.5cm}>{\raggedright}p{4.0cm}>{\raggedright}p{5.6cm}>{\raggedright}p{2.3cm}>{\raggedright\arraybackslash}p{1.7cm}}
\toprule
No. & Claim & Experiment and what was measured & Source & Status \\
\midrule\endfirsthead
\toprule
No. & Claim & Experiment and what was measured & Source & Status \\
\midrule\endhead
1 & The brain's \nt{ACET} neurons are few, gathered in a few groups, and their outputs spread throughout the cortex: a broadcast channel & Antibody staining for the acetylcholine-synthesizing enzyme and retrograde labeling in the rat: the cortex, hippocampus, amygdala, olfactory bulb and thalamus receive acetylcholine mainly from six compact cell groups of the basal forebrain and brainstem & \cite{Mesulam1983} & measured \\
2 & The acetylcholine level in the cortex is high in waking and low in slow-wave sleep & Microdialysis in the cortex and hippocampus of freely moving cats with EEG recording: acetylcholine release is $\sim100\%$ higher in quiet waking and $\sim175\%$ higher in active waking than in slow-wave sleep; in REM sleep, cortical release is as in waking & \cite{Marrosu1995,Hasselmo1999} & measured \\
3 & The meaning of the signal is set by the receptor: acetylcholine has fast channel receptors and slow ones that trigger reactions in the cell (Proposition~\ref{prop:receptor}) & Review of measurements: the action of acetylcholine on excitability, transmission and plasticity depends on the release site, the receptor subtype (nicotinic are ion channels, muscarinic act through G proteins) and the recipient cell & \cite{Picciotto2012} & measured \\
4 & First action: weaken the internal connections while hardly touching external inputs (factor $1-a$ in~\eqref{eq:acetnet}) & Rat olfactory cortex slices: at $100$\,\textmu M acetylcholine agonists, potentials of internal connections (layer Ib) fall by more than $60\%$, those of the external input (layer Ia) by less than $15\%$, in the same cell; the suppression is presynaptic. Hippocampal slices (CA1): $89\%$ versus $40\%$ for the internal and input pathways & \cite{HasselmoBower1992,HasselmoSchnell1994} & measured \\
5 & Second action: strengthen the input (factor $\frac{1+a}{2}$) & Neocortical slices: nicotinic receptors strengthen synapses of the thalamic input and do not act on intracortical ones; muscarinic receptors weaken both types. Acetylcholine raises neuronal excitability (review) & \cite{Gil1997,Hasselmo2006} & measured \\
6 & Third action: make weight changes easier (rate $\eta a$) & Rat: electrical stimulation of the nucleus basalis (the source of cortical acetylcholine) paired with a tone produces, in an adult animal, a strong reorganization of the auditory cortex map toward the presented sound & \cite{KilgardMerzenich1998,Hasselmo2006} & measured \\
7 & The Hebb rule for sparse patterns in the second equation of~\eqref{eq:acetnet} gives an associative memory that completes a pattern from a part & Capacity calculation for a network with sparse patterns and a covariance rule: at a small fraction $p$ of active neurons, the network stores considerably more patterns than at $p=1/2$ & \cite{Tsodyks1988} & theory \\
8 & Writing at low $a$ corrupts memory: the network writes down what it recalled, not what it saw; at $a=0.1$ the corruption is no weaker than at $a=0$ & \texttt{models/\allowbreak acet\_memory.py}: $200$ neurons, five patterns of $20$, a new pattern shares $10$ neurons with one of them, $10$ runs. At $a=0$ the new pattern is not memorized, and the old one drags along $5.9$ foreign neurons out of $10$; at $a=0.1$ the new pattern is recalled ($0.97$), but the two merge: the new one drags along $7.3$ foreign neurons, the old one $7.9$; from $a=0.2$ on, fewer than one & derived in the chapter & model \\
9 & Proposition~\ref{prop:writeread}: there is no level $a$ at which writing and reading are equally good (Fig.~\ref{fig:acetsim}) & The same script, learning rate $\eta a$: a new pattern is memorized in full after one presentation for $a\ge0.9$, to $0.8$ at $a=0.75$, not at all for $a\le0.5$. Reading from half a pattern: $1.0$ for $a\le0.2$, $0.96$ at $a=0.3$, $0.04$ at $a=0.4$ & derived in the chapter & model \\
10 & In the brain a single broadcast wire, \nt{ACET}, switches between writing and reading & The idea comes from models built on slice measurements. The most direct experiment is in humans: the muscarinic blocker scopolamine impairs learning of new word pairs, most of all those overlapping with old ones, but not recall of pairs already learned. There is no direct measurement of two network modes & \cite{HasselmoAndersonBower1992,Atri2004} & hypothesis \\
11 & The filter~\eqref{eq:kalman} is optimal; the input weight and the learning rate are one quantity, $P/R$; in steady state $P=\sqrt{qR}$ and the memory is $\sqrt{R/q}$ (Proposition~\ref{prop:kalman}) & Requires no experiment: the optimality of the Kalman--Bucy filter is a classical result of estimation theory; the steady state is derived in the chapter & \cite{KalmanBucy1961} & theory \\
12 & \nt{ACET} tells the network a quantity like $P$: the expected uncertainty & Proposed in a probabilistic model of attention. There is no direct measurement showing that the acetylcholine level tracks the uncertainty of an estimate & \cite{YuDayan2005} & hypothesis \\
13 & Some \nt{ACET} neurons respond to reward and punishment within twenty milliseconds or so, more strongly the more unexpected the reinforcement & Mouse, optogenetically identified cholinergic neurons of the basal forebrain: response to reward and punishment after $18\pm3$\,ms; the size grows with the surprise of the reinforcement; responses are similar in neurons of two nuclei & \cite{Hangya2015} & measured \\
14 & \nt{NORA} notices an unexpected change in the world, and \nt{ACET} keeps the network in write mode & The division of labor is proposed in a model. Indirectly: in humans, pupil diameter, linked to \nt{NORA}, tracks changes and the learning rate. There is no direct test of the \nt{NORA}--\nt{ACET} pair & \cite{YuDayan2005,Nassar2012} & hypothesis \\
15 & In slow-wave sleep the network, in read mode, replays what was written during the day, and these replays rewrite it into long-term memory & Recordings of rat hippocampal ensembles: cells that worked together during the day fire together more often in subsequent slow-wave sleep, and in the same order (measured). For rewriting: suppressing hippocampal ripples after learning impairs spatial memory, and lengthening spontaneous ripples improves it; that the cause is a low \nt{ACET} level has not been proven & \cite{WilsonMcNaughton1994,Skaggs1996,Girardeau2009,FernandezRuiz2019,Hasselmo1999} & hypothesis \\
\bottomrule
\end{longtable}}

\section{Chapter~\ref{sec:synthesis}. The Whole Machine}

This summary chapter brings in no new experiments: each row rests on the chapter where the claim is worked out and on the same sources. The \nt{GLUT} bus and flow control through \nt{GABA} are firmly established, while the roles of the broadcast channels and their joint operation are mostly hypothesis.

{\footnotesize
\setlength{\tabcolsep}{3pt}\renewcommand{\arraystretch}{1.15}
\begin{longtable}{>{\raggedright}p{0.5cm}>{\raggedright}p{4.0cm}>{\raggedright}p{5.6cm}>{\raggedright}p{2.3cm}>{\raggedright\arraybackslash}p{1.7cm}}
\toprule
No. & Claim & Experiment and what was measured & Source & Status \\
\midrule\endfirsthead
\toprule
No. & Claim & Experiment and what was measured & Source & Status \\
\midrule\endhead
1 & There are only four kinds of control signal for $8.6\cdot10^{10}$ neurons~\eqref{eq:machine} & Counting cell nuclei in brain homogenate (isotropic fractionator): an adult man has $86.1\pm8.1$ billion neurons & \cite{Azevedo2009} & measured \\
2 & Proposition~\ref{prop:architecture}, the first two layers: data travel along the \nt{GLUT} address bus, the sign of a connection is set by the sender, and \nt{GABA} neurons steer and normalize the flow (Chapters~\ref{sec:neurontypes}, \ref{sec:gaba}) & One main transmitter per neuron (review of measurements); feedforward inhibition narrows the window in which a pyramidal neuron sums its inputs; division by the total activity has been measured in many areas & \cite{Strata1999,Pouille2001,Carandini2012} & measured \\
3 & The elements are unreliable: a synapse loses most spikes (Table~\ref{tab:computer}, Chapter~\ref{sec:code}) & Hippocampal slices, minimal stimulation: more than half of the spikes evoke no response in CA1; threshold failures and axonal conduction failures are ruled out, and the cause is probabilistic transmitter release & \cite{AllenStevens1994} & measured \\
4 & The brain runs on $\sim20$ watts, on average about one spike per second per neuron (Chapter~\ref{sec:code}); engineers have not yet built such a machine & The estimate~\eqref{eq:energy} from measured costs: about $10^9$ ATP molecules per spike, including synaptic work (rodent cortex); a detailed estimate for the cortex gives an even lower rate. The state of the art is from a review & \cite{Attwell2001,Lennie2003,MehonicKenyon2022} & theory \\
5 & Learning at most synapses is the product of a local trace and one shared number $\delta$ (the second equation of~\eqref{eq:machine}, Chapter~\ref{sec:dofa}) & Monkey \nt{DOPA} neurons respond to the reward-prediction error; in striatal slices dopamine enlarges a spine only if it arrives $0.3$--$2$\,s after the glutamate input, so the trace has been measured & \cite{Schultz1997,Yagishita2014} & measured \\
6 & A separate error wire to each output exists only in the circuit for learning movements (Chapter~\ref{sec:motorlearning}) & Cerebellum: coincidence of the climbing-fiber signal with an input weakens that input; in the monkey, complex spikes of Purkinje cells carry the direction of the movement error and cause a correction & \cite{Ito2001,Herzfeld2018} & measured \\
7 & Each broadcast channel is a bundle of several lines (Proposition~\ref{prop:bundle}) & For \nt{DOPA}: in one task different neurons carry reward, movement, and stimulus features; the fast error and the slow dopamine level are distinct. For \nt{SERO}, \nt{NORA}, and \nt{ACET} such a breakdown is less well studied & \cite{Engelhard2019,Hamid2016,Mohebi2019} & measured \\
8 & \nt{SERO}: a level regulator and, by hypothesis, the horizon $\tau_\gamma$ (Chapter~\ref{sec:serocomputer}) & Self-inhibition: agonists of their own 5-HT$_{1A}$ receptors inhibit the serotonin neurons of the raphe (measured). The horizon was proposed in theory; the strongest experiment is that optogenetic activation of these neurons in the mouse lengthens waiting for a delayed reward & \cite{Sprouse1987,Doya2002,Miyazaki2014} & hypothesis \\
9 & \nt{NORA}: the gain $g$ chooses between searching and deciding (Chapter~\ref{sec:nora}) & Rise in signal-to-noise ratio: in the awake monkey, noradrenaline in the auditory cortex suppresses background activity more strongly than the response to conspecific calls (measured); multiplicative gain is a model; the role in choosing between search and decision is theory & \cite{Foote1975NE,ServanSchreiber1990,AstonJones2005} & hypothesis \\
10 & \nt{ACET}: switches between writing and reading (Chapter~\ref{sec:acet}) & Three actions (weakening internal connections, amplifying the input, facilitating plasticity) are measured; mode switching is indirectly supported by an experiment with scopolamine in humans & \cite{HasselmoBower1992,Gil1997,Atri2004} & hypothesis \\
11 & Formula~\eqref{eq:machine} describes the whole machine & A summary of the models of Chapters~\ref{sec:glutcomputer}--\ref{sec:acet}, each part resting on its own chapter; as a whole it has not been tested experimentally & derived in the chapter & hypothesis \\
12 & The knobs work in concert without central control: error, surprise, writing, return (Fig.~\ref{fig:timeline}) & No experiment: the sketch is qualitative. Individual links are the theory of the division of labor between \nt{NORA} and \nt{ACET} and the link between pupil size and learning rate in humans & \cite{YuDayan2005,Nassar2012} & hypothesis \\
13 & Learning from one shared signal is slower the more neurons affect the result (Proposition~\ref{prop:broadcastcost}) & Calculation of the learning curves of a linear network: learning by output perturbation is slower than gradient descent by a factor equal to the number of outputs & \cite{Werfel2005} & theory \\
14 & How the cortex tunes multilayer circuits is unknown; the candidates are bursts of spikes, computation in the neuron's branches, and self-supervised prediction & Bursts as carriers of error: model; only a $5$--$8$-layer network could reproduce the response of a detailed model of a cortical neuron: model; calcium spikes in the branches of human cortical neurons that give exclusive OR: measured in slices; self-supervised learning: a review of the evidence & \cite{Payeur2021,Beniaguev2021,Gidon2020,Keller2018} & hypothesis \\
\bottomrule
\end{longtable}}

\section{Chapter~\ref{sec:sensors}. The Machine's Inputs}

The structure of the sensors has been measured directly, by recording currents from single cells and vibrations of the cochlea and by counting retinal cells. The sensitivity limit and the shot-noise formula follow from physics, while the claim that the sensors' codes are tuned for the greatest information is a hypothesis with strong support.

{\footnotesize
\setlength{\tabcolsep}{3pt}\renewcommand{\arraystretch}{1.15}
\begin{longtable}{>{\raggedright}p{0.5cm}>{\raggedright}p{4.0cm}>{\raggedright}p{5.6cm}>{\raggedright}p{2.3cm}>{\raggedright\arraybackslash}p{1.7cm}}
\toprule
No. & Claim & Experiment and what was measured & Source & Status \\
\midrule\endfirsthead
\toprule
No. & Claim & Experiment and what was measured & Source & Status \\
\midrule\endhead
1 & A sensor is a transducer: the receptor gives a current, and the output of the sensory cells of the eye and ear is glutamate onto the \nt{GLUT} bus; the first cells produce no spikes (Fig.~\ref{fig:transducer}) & Turtle rods and cones release an excitatory amino acid, detected by glutamate channels on an excised membrane patch. Inner hair cells of the rat cochlea: currents in the nerve-fiber endings flow through glutamate AMPA receptors, and the release rate grows with smooth depolarization of the cell up to $150$\,Hz & \cite{CopenhagenJahr1989,GlowatzkiFuchs2002} & measured \\
2 & A light sensor in the dark responds to a single photon with a current of about a picoampere & Suction electrode on the outer segment of a toad rod: responses to dim flashes are quantized, an event of $\approx1$\,pA ($5\%$ of the dark current) with a spread of $0.2$\,pA, efficiency $0.5$. Humans report the arrival of a single photon more often than chance & \cite{Baylor1979,Tinsley2016} & measured \\
3 & A sound sensor detects a displacement of less than a nanometer & M\"ossbauer technique, basilar membrane of the guinea-pig cochlea at the $16$--$19$\,kHz place: at the auditory-nerve response threshold the membrane velocity is about $0.04$\,mm/s, i.e., a displacement of $\approx0.35$\,nm (our conversion). The membrane was measured, not the sensor's bundle & \cite{Sellick1982,Hudspeth1989} & measured \\
4 & Precision in the dark is limited by photon shot noise~\eqref{eq:photonnoise}; in dim light integration is longer & The formula follows from Poisson statistics. In a rod the noise in dim light matches the sum of independent quantal events; against a bright background the response to a flash is smaller and shorter. The ``tenths of a second'' figure is not checked separately here & \cite{Baylor1979} & theory \\
5 & The input range is ten orders of magnitude for light and twelve for sound, while the spike rate spans less than three & For sound: the basilar-membrane response at the characteristic frequency grows with a slope as low as $0.2$\,dB/dB over $40$--$80$\,dB, and compression is seen even above $100$\,dB. The order-of-magnitude figures themselves are standard estimates, unsourced in the chapter & \cite{Ruggero1997} & measured \\
6 & The sensor response is~\eqref{eq:nakarushton}, and $\sigma$ follows the mean brightness, shifting the curve along the logarithmic axis (Fig.~\ref{fig:adaptation}) & The formula was first fitted to S-potentials of the fish retina. Turtle cones: responses obey $V=I V_m/(I+\sigma)$ and span $\sim3.5$ orders of magnitude of brightness; after $2$--$3$ minutes of background the curve shifts horizontally, keeping its width. The link to division by total activity is from a review & \cite{NakaRushton1966,NormannPerlman1979,Carandini2012} & measured \\
7 & Sensors transmit changes, and their codes are tuned for the greatest information per spike (Proposition~\ref{prop:sensors}) & The principle was proposed theoretically. The strongest evidence: in the fly, the ``contrast $\to$ response'' curve of first-layer neurons matches the cumulative distribution of contrasts in natural scenes, which yields the greatest information & \cite{Barlow1961,Laughlin1981} & hypothesis \\
8 & On and off sensors form a two-wire code; the event camera reproduces it in silicon & A review of experiments: the on and off channels of the retina separate already at the first synapse and can be switched off separately. Camera: $128\times128$ pixels without frames, $120$\,dB range, $15$\,\textmu s latency & \cite{Schiller1992,Lichtsteiner2008,Gallego2022} & measured \\
9 & The eye has $\sim10^8$ light sensors and $\sim10^6$ output neurons: compression by $\sim100$ & Counts on whole-mount human retinas: on average $92$ million rods and $4.6$ million cones; $0.7$ to $1.5$ million output (ganglion) cells in different eyes & \cite{Curcio1990,CurcioAllen1990} & measured \\
10 & The mouse eye has more than thirty output channels & Mouse: two-photon calcium imaging of more than $11\,000$ output cells and clustering of the responses gives clearly more than $30$ functional types. The number for humans has not been measured & \cite{Baden2016} & measured \\
11 & The guinea-pig eye transmits $\sim875\,000$ bits/s; scaling to the human eye gives $\sim10^7$ bits/s & Guinea pig, multielectrode array, motion in natural scenes: $6$--$13$ bits/s per cell, $\sim875\,000$ bits/s for $\sim10^5$ output cells. For humans ($\sim10^6$ cells), $\sim10^7$ bits/s is the authors' extrapolation & \cite{Koch2006} & measured \\
12 & The ear is a bank of bandpass filters with a nearly logarithmic frequency map (Fig.~\ref{fig:filterbank}) & The place of maximum basilar-membrane vibration depends on frequency; the frequency--place function is nearly exponential and a single form describes the cochleas of humans and living animals & \cite{Greenwood1990,Robles2001} & measured \\
13 & The resonators are active: some sensors amplify weak vibrations thousands of times in amplitude ($66$--$76$\,dB), and the gain falls as loudness grows & Laser vibrometry at the base of the chinchilla cochlea: for weak tones at the characteristic frequency the gain is $66$--$76$\,dB relative to the stapes; death lowers sensitivity by $60$--$81$\,dB (a factor of $10^3$--$10^4$ in amplitude) and removes the compression. The source of the force is the outer hair cells & \cite{Ruggero1997,Robles2001} & measured \\
14 & The ear's amplifier sits at the edge of self-oscillation & Calculation: near a Hopf bifurcation point, range compression, sharp tuning, and combination tones, all the known nonlinearities of hearing, are unavoidable. That the cochlea is tuned in exactly this way has not been proved directly & \cite{Eguiluz2000} & hypothesis \\
15 & At low frequencies spikes are phase-locked to the sound (in the guinea pig, locking weakens from $\sim600$\,Hz and is noticeable up to $\sim3.5$\,kHz), and direction is found from them; at high frequencies only the place code remains & Guinea-pig auditory nerve: phase locking begins to fall at about $600$\,Hz and disappears above $3.5$\,kHz; in the cat and the monkey the limit is about an octave higher. The interaural time difference is from a review & \cite{PalmerRussell1986,Grothe2010} & measured \\
16 & The other sensors (Table~\ref{tab:sensors}): smell has $\sim400$ receptor types, one per neuron, and a combinatorial code; taste works partly through ATP and serotonin; touch has fast- and slow-adapting sensors; warm and cold are separate & Mouse: one receptor recognizes many odors, one odor activates many receptors, and different odors activate different sets. Without ATP receptors the taste nerve does not respond; taste buds release serotonin. Single-fiber recordings from the skin of the human hand: four types by adaptation speed. The cold channel is distinct from the heat channels & \cite{Mombaerts2004,Malnic1999,Finger2005,Huang2005,JohanssonVallbo1979,McKemy2002} & measured \\
\bottomrule
\end{longtable}}

\section{Chapter~\ref{sec:recognition}. Recognition}

The speed of recognition, the structure of the first stages, the linear readout of the response, and learning from the continuity of time are measured in humans and monkeys; the reduction of the whole chain to two operations and learning from video are tested on the book's models, and the explanations of illusions range from well founded to disputed.

{\footnotesize
\setlength{\tabcolsep}{3pt}\renewcommand{\arraystretch}{1.15}
\begin{longtable}{>{\raggedright}p{0.5cm}>{\raggedright}p{4.0cm}>{\raggedright}p{5.6cm}>{\raggedright}p{2.3cm}>{\raggedright\arraybackslash}p{1.7cm}}
\toprule
No. & Claim & Experiment and what was measured & Source & Status \\
\midrule\endfirsthead
\toprule
No. & Claim & Experiment and what was measured & Source & Status \\
\midrule\endhead
1 & Pixel-by-pixel template matching under shift does no better than a coin; a pair of stages~\eqref{eq:scstage} recognizes the shape at any position & \texttt{models/\allowbreak recognition\_models.py}, a ``retina'' of $24$ points, $4000$ trials: template, $0.49$, $0.51$, $0.52$ at flipped-point fractions $0$, $0.1$, $0.2$; the $S$ and $C$ pair, $1.00$, $0.86$, $0.70$ & derived in the chapter & model \\
2 & An animal in a picture is recognized in $\sim150$\,ms, a single pass through about ten stages of $10$--$20$\,ms each & Humans, ``animal or not'' task on photographs shown for $20$\,ms: evoked potentials for the two answers diverge after $\sim150$\,ms. Monkeys: the latency of the first response grows from stage to stage (V1 earliest, then V2 and V4). The number of stages and ``zero or one spike per stage'' are inferred from these numbers & \cite{Thorpe1996,Schmolesky1998} & measured \\
3 & Stage $S$ is AND over parts, stage $C$ is the maximum over positions (Proposition~\ref{prop:conveyor}) & Cat, primary visual cortex: simple cells respond to an edge of a given orientation at a given place, complex cells to the same orientation over a larger region. Intracellular recording from complex cells: for many cells the response to two bars is close to the larger of the two responses, on average closest to a max operation. The maximum via mutual inhibition is Proposition~\ref{prop:wta}; division by the total activity, a review & \cite{HubelWiesel1962,Lampl2004,Carandini2012} & measured \\
4 & Further up the chain the combinations are larger and more complex, and the response depends less and less on position & Monkeys, stimuli simplified down to the ``critical feature'': cells of V4 and posterior inferotemporal cortex respond to complex combinations of shape, color, and texture with small fields; in anterior inferotemporal cortex the fields are larger. Alternating $S$ and $C$ in a model reproduces this picture & \cite{KobatakeTanaka1994,Riesenhuber1999} & measured \\
5 & Convolutional networks repeat this alternation and predict the responses of the upper stages best; the object is read out linearly from the rates of a group of neurons & A network trained for recognition, with no fitting to the brain: the top layer predicts the responses of monkey inferotemporal neurons, the middle layers the responses of V4. A linear classifier on the activity of $\sim100$ randomly chosen cells, in windows from $12.5$\,ms, recognizes object and category across positions and sizes & \cite{Fukushima1980,Yamins2014,Hung2005} & measured \\
6 & The Hebb rule with a trace~\eqref{eq:traceinvariance} learns invariance from unlabeled video if the trace is at least $\sim20$\,ms & The idea of slow features and the trace rule is theory. \texttt{models/\allowbreak recognition\_models.py}, two $C$ neurons, $16$ inputs, $10$ runs, a $0.3$\,s pause between objects. At $\tau_{\text{tr}}=5$\,ms the match with the shape is $0.52$ on average, at $10$\,ms $0.56$; at $20$, $50$, and $200$\,ms, $1.00$ in all runs (at $30$\,ms one run in ten errs) & \cite{Foldiak1991,WiskottSejnowski2002} & model \\
7 & In the brain, invariance is learned from the continuity of time & Monkeys: the object was covertly swapped during an eye jump to a particular place in the field; the position invariance of inferotemporal neurons changed in line with the swap, significantly after just an hour. That this is the main rule of visual learning is a hypothesis & \cite{LiDiCarlo2008} & measured \\
8 & Motion: direction detectors, then the same AND/OR pair over large regions & Direction detectors in the rabbit retina; in monkey area MT all $110$ recorded neurons are direction selective, and the response to motion is stronger than in V1 & \cite{Barlow1965,Albright1984} & measured \\
9 & The upper stages send a prediction down, and the error goes up & The model explains extra-classical receptive field effects in primary visual cortex; individual observations agree (review), but it is not proven as the general design of the pipeline & \cite{RaoBallard1999,Keller2018} & hypothesis \\
10 & Hard pictures require feedback and several rounds & Monkeys: for ``hard'' pictures the decision in inferotemporal cortex forms $\sim30$\,ms later; these late responses are poorly predicted by a feedforward network and better by networks with feedback or by deeper ones & \cite{Kar2019} & measured \\
11 & An invisible pixel forgery fools a network; human vision is far more robust, but not immune & In networks: a small, specially chosen perturbation changes the answer; the ``panda $\to$ gibbon'' example is from the second paper. In humans, with brief presentation, forgeries that transfer well between networks shift the choice & \cite{Szegedy2014,Goodfellow2015,Elsayed2018} & measured \\
12 & Illusions of expectation: an unreliable input yields to expectation; faint things move ``more slowly,'' the dress is seen differently (Section~\ref{sec:illusions}) & Lower-contrast gratings appear to move more slowly. A survey of $1401$ people: the dress is seen in three stable variants, and a lighting cue switches the color; in $\sim13\,000$ observers the assumption about lighting decides. A Bayesian model with an expectation of slow motion explains a range of motion illusions & \cite{Thompson1982,LaferSousa2015,Wallisch2017,Weiss2002,Purves2011} & hypothesis \\
13 & Gain control: the waterfall, tilt, and contrast; the Hermann grid is not inhibition from neighbors & Direction detectors in the rabbit retina lower their rate under prolonged motion and fall below baseline after it stops. The tilt illusion is reproduced by a model with division by the activity of neighbors. Grid modifications that do not change the response of retinal output neurons abolish the spots, so the neighbor-inhibition explanation is rejected & \cite{BarlowHill1963,Schwartz2009,SchillerCarvey2005} & measured \\
14 & Delay compensation: a moving object appears ahead of a flash & The illusion is measured. Psychophysics contradicts both forward extrapolation of motion and a latency difference; a retrospective explanation has been proposed, based on events $\sim80$\,ms after the flash. The debate is not settled & \cite{Nijhawan1994,EaglemanSejnowski2000} & hypothesis \\
15 & Still ``snakes'' rotate: the latency difference~\eqref{eq:latency} is read by direction detectors & The illusion works without color too; pairs of neighboring elements produce it on their own; direction-selective neurons in monkey cortex give a directional response to these still pairs. In humans, rotation is triggered by microsaccades and blinks & \cite{Conway2005,OteroMillan2012} & measured \\
16 & Ambiguous pictures: mutual inhibition, fatigue, and noise; switches are caused mainly by noise & A model of competing neuron groups: switches in it are caused by \emph{noise}, and without noise there are none; it explains why the switching rate grows with stimulus strength. Fatigue plays a smaller role here & \cite{MorenoBote2007} & hypothesis \\
17 & Some illusions follow from the task the machine learns & A network trained to predict the next video frame ``sees'' rotation of the still rings and does not see it in control pictures; networks for image denoising and sharpening fall for color and brightness illusions, as humans do & \cite{Watanabe2018,GomezVilla2020} & measured \\
\bottomrule
\end{longtable}}

\section{Chapter~\ref{sec:muscles}. Output to muscles}

The design of the output (motor units, the synaptic safety factor, recruitment by size, the muscle pair, the stretch loop, and rhythm generators) is measured directly; the stability condition for a loop with delay is a theorem; the internal model of the body is a well-founded hypothesis; the numbers in the figures are computed with the book's model.

{\footnotesize
\setlength{\tabcolsep}{3pt}\renewcommand{\arraystretch}{1.15}
\begin{longtable}{>{\raggedright}p{0.5cm}>{\raggedright}p{4.0cm}>{\raggedright}p{5.6cm}>{\raggedright}p{2.3cm}>{\raggedright\arraybackslash}p{1.7cm}}
\toprule
No. & Claim & Experiment and what was measured & Source & Status \\
\midrule\endfirsthead
\toprule
No. & Claim & Experiment and what was measured & Source & Status \\
\midrule\endhead
1 & Motor unit: one \nt{ACET} neuron controls a group of fibers, from several hundred to a couple of thousand; in muscles for fine adjustment the units are smaller & Human muscles, counts of motor axons and muscle fibers: on average about $340$ fibers per neuron in an interosseous muscle of the hand, about $410$ in the brachioradialis, about $1930$ in the medial gastrocnemius. The chapter gives no exact number for the smallest units. & \cite{Feinstein1955mu} & measured \\
2 & The synapse on a muscle releases several times more transmitter than is needed to make the fiber fire & Review of measurements at neuromuscular synapses: the safety factor (ratio of released transmitter to threshold) in adult mammals is usually $3$--$5$; in humans the margin rests more on the folds of the receiving membrane than on the size of the release. & \cite{WoodSlater2001} & measured \\
3 & The twitch~\eqref{eq:twitch} rises over a time $T_c$ of order $50\ms$ & Single motor units of the human hand during voluntary effort, force averaged over the unit's spikes: rise time $30$--$100\ms$, under $70\ms$ for $80\,\%$ of units. The function $(t/T_c)e^{1-t/T_c}$ is an approximation from a motor-unit pool model. & \cite{MilnerBrown1973rec, Fuglevand1993} & measured \\
4 & Sum of twitches~\eqref{eq:forcesum}: at $5$, $15$, and $40$ spikes/s the force is $0.2$--$1.1F_1$, $1.8$--$2.2F_1$, and about $5.4F_1$ (Fig.~\ref{fig:tetanus}) & Computed with \texttt{models/\allowbreak muscle\_\allowbreak models.py}, $T_c=50\ms$: $0.21$--$1.10$, $1.76$--$2.19$, and $5.28$--$5.49\,F_1$ over the last $0.3\,\text{s}$. Linear summation is a simplification: the model has no saturation of a real muscle. & \texttt{models/\allowbreak muscle\_\allowbreak models.py} & model \\
5 & Units are recruited by size; the weakest are a hundred times weaker than the strongest & Cat ventral roots: with increasing reflex input, the neurons with small spikes in the root, that is, the small ones, fire first. Human hand interosseous muscle: unit twitch forces range from $0.1$ to $10$~g and grow almost linearly with the force at which the unit is recruited. & \cite{Henneman1957, MilnerBrown1973rec} & measured \\
6 & The force step is proportional to the force, $\Delta F/F\approx q-1$~\eqref{eq:sizeprinciple}: floating point & Derived in the chapter from the geometric series of forces. Consistent with experiment: at large efforts, far fewer new units are recruited for the same force increment. & derived in the chapter; \cite{MilnerBrown1973rec} & theory \\
7 & The scatter of force grows in proportion to the force itself & Voluntary isometric effort of a thumb muscle: the standard deviation of force grows linearly with mean force; for the same effort evoked by electrical stimulation, the scatter is constant. Pool model: the linear growth comes from recruiting units in order of force. & \cite{Jones2002sdn} & measured \\
8 & The smoothness of movements is a consequence of signal-dependent noise & Model: the trajectory that minimizes end-point scatter under such noise reproduces measured eye and arm trajectories. There is no direct experiment that separates this cause from others. & \cite{HarrisWolpert1998} & hypothesis \\
9 & The difference of the forces of a muscle pair sets the torque, the sum sets the stiffness~\eqref{eq:antagonists} & Theory of impedance control by co-contraction. Human arm: the stiffness of each joint, measured with small perturbations, is roughly proportional to the torque at it; different co-contraction ratios change the shape and orientation of the hand stiffness ellipse. & \cite{Hogan1984, GomiOsu1998stiff} & measured \\
10 & The stretch sensor excites its own \nt{ACET} neuron and, through an inhibitory intermediary, switches off the opponent (Fig.~\ref{fig:antagonists}) & Cat spinal cord: a single intermediary excited by stretch-sensor fibers produces monosynaptic inhibitory potentials in the \nt{ACET} neurons of the opposing muscle, synaptic delay $0.28$--$0.42\ms$. The intermediary's transmitter (glycine) was not identified in this experiment. & \cite{Jankowska1972Ia} & measured \\
11 & Loop delay: spinal $25$--$30\ms$, through the cortex one and a half to three times longer, through the eye $100$--$200\ms$ & Human arm, a perturbation of the joint: short-latency muscle response at $20$--$45\ms$, long-latency at $45$--$105\ms$, voluntary at $120$--$180\ms$. A shift in the seen position of the finger during a movement: correction after $160\ms$ on average. & \cite{Pruszynski2008rapid, SaundersKnill2003vis} & measured \\
12 & $\dd x/\dd t=-Gx(t-d)$ is stable for $Gd<\pi/2$~\eqref{eq:delaystable}; at the boundary the frequency is $1/(4d)$ & Theorem on the roots of a delay equation. Check with \texttt{models/\allowbreak muscle\_\allowbreak models.py}, $d=30\ms$: by $3\,\text{s}$ the amplitude is $3\cdot10^{-6}$ at $G=40$, $0.13$ at $G=50$, $38$ at $G=55$ per second (boundary $52$). & \cite{Hayes1950} & theory \\
13 & Physiological tremor at $8$--$12$~Hz; the stretch loop is one of its sources & Review of measurements: fine tremor around $10$~Hz is present in healthy people; its origin is mixed (central oscillations, properties of unit discharges, mechanical resonance, and resonance of the reflex loop), and different sources dominate under different conditions. & \cite{McAuleyMarsden2000} & hypothesis \\
14 & The brain predicts the result of a command from an internal model of the body & A person holds a load in a pinch grip and moves the arm: under inertial, viscous, and elastic loads, grip force changes simultaneously with the load, not with the loop delay, so the load is predicted. A review collects such experiments; the model itself has not been observed directly in neurons. & \cite{FlanaganWing1997grip, WolpertGhahramani2000} & hypothesis \\
15 & The rhythm of walking, breathing, and chewing is generated by local networks under constant input & Review of experiments: an isolated nervous system (spinal cord, ganglia) produces a rhythmic motor pattern without rhythmic input and without feedback from muscles. & \cite{MarderBucher2001} & measured \\
16 & Mutual inhibition with fatigue~\eqref{eq:halfcentre} gives a rhythm with period $0.84\,\text{s}\approx2\tau_a$ (Fig.~\ref{fig:halfcentre}) & Theorem: a network with mutual inhibition and adaptation that has no stable state under constant input must oscillate. Computed with \texttt{models/\allowbreak muscle\_\allowbreak models.py} ($I=1$, $\beta=2$, $b=2$, $\tau=20\ms$, $\tau_a=0.4\,\text{s}$): period $0.84\,\text{s}$. That real generators are built exactly this way has not been shown. & \cite{Matsuoka1985osc} & model \\
\bottomrule
\end{longtable}}

\section{Chapter~\ref{sec:pain}. Pain}

Pain sensors, the two wires, the gate in the spinal cord, the descending volume knob, sensitization, and the capture of attention have been measured directly; the gate and sensitization formulas are simplified models of the book; the rule by which the machine chooses the volume of the alarm is a hypothesis.

{\footnotesize
\setlength{\tabcolsep}{3pt}\renewcommand{\arraystretch}{1.15}
\begin{longtable}{>{\raggedright}p{0.5cm}>{\raggedright}p{4.0cm}>{\raggedright}p{5.6cm}>{\raggedright}p{2.3cm}>{\raggedright\arraybackslash}p{1.7cm}}
\toprule
No. & Claim & Experiment and what was measured & Source & Status \\
\midrule\endfirsthead
\toprule
No. & Claim & Experiment and what was measured & Source & Status \\
\midrule\endhead
1 & High threshold: heat thresholds of different pain sensors range from $41^\circ$ to $46$--$48^\circ$ & Single-fiber recordings. Monkey skin: median threshold of the slow (C) sensors $41^\circ$, of fast type~II sensors $46^\circ$, of type~I above $53^\circ$. Human skin: C sensors that also respond to pressure $40.7^\circ$, those that do not respond to pressure $48^\circ$. & \cite{Treede1995heat, Weidner1999cfib} & measured \\
2 & Pain sensors adapt far less than touch sensors and become more sensitive after mild injury; a weakening after strong heating is measured in one type & A mild skin burn ($50^\circ$, $100\,\text{s}$): after $5$--$10$~min the pain threshold in humans and the threshold of C sensors in the monkey are lowered by $2$--$6^\circ$; other skin sensors show no such shift. In fast type~II sensors the response after strong heating is suppressed. The comparison of adaptation with touch sensors is a general estimate, not backed here by a separate experiment. & \cite{LaMotte1982hyper, Treede1995heat} & measured \\
3 & Two wires: fast $5$--$30$~m/s, slow $0.5$--$2$~m/s; arrival time $t=L/v$~\eqref{eq:paintime} & Conduction velocity: fast pain sensors of the monkey average $15$ and $25$~m/s (two types); human C sensors $0.8$--$1.0$~m/s. For $L=1$~m this gives tens of milliseconds and about a second. & \cite{Treede1995heat, Weidner1999cfib} & measured \\
4 & Pain arrives twice: first fast and precise, then dull and lasting & Reaction time to heating of the human forearm: from $1100$ to $400\ms$ as the temperature rises; strong heating of hairy skin gives double pain, of the palm does not. First pain can be carried only by fibers faster than $6$~m/s; by latency it corresponds to fast type~II sensors, which are absent from the palm. The division of roles (``where'' and ``how bad'') is an interpretation. & \cite{CampbellLaMotte1983first, Treede1995heat} & measured \\
5 & The gate: an output neuron of the spinal cord receives pain and touch, and an inhibitory interneuron is excited by touch (Fig.~\ref{fig:gate}) & The gate scheme was proposed as a theory. Mice, genetic labeling and ablation of neuron groups: excitatory somatostatin neurons are needed for mechanical pain; inhibitory dynorphin neurons keep input from touch sensors from reaching them, and without these neurons a light touch causes pain. & \cite{MelzackWall1965, Duan2014} & measured \\
6 & Touch dampens pain & Humans, laser pain pulses on the arm: simultaneous touch lowers the pain rating and the ability to discriminate pulse energy; the effect weakens with the distance between the touch point and the pain point within one skin segment. & \cite{Mancini2014touch} & measured \\
7 & Assumption of gate theory: strong pain itself switches off the inhibitory interneuron, and the gate swings open & Marked in the chapter as an assumption of the original theory. The mouse work describes how the gate keeps touch out of the pain channel; switching off of the interneuron by pain is not shown in it. & \cite{MelzackWall1965} & hypothesis \\
8 & The gate~\eqref{eq:gate}: threshold $0.18$ without touch, $0.86$ with touch, $0.11$ at $g=2$; at $\nu=1$ touch lowers the output from $1$ to $0.25$, at $\nu=2$ it has no effect; touch alone does not get through (Fig.~\ref{fig:gatecurves}) & A calculation with \texttt{models/\allowbreak pain\_\allowbreak gate.py} using the weights of the text reproduces all these numbers. & \texttt{models/\allowbreak pain\_\allowbreak gate.py} & model \\
9 & Descending pathways from the brainstem change the gain of the gate in both directions & Rat, medulla, recordings during tail heating: some neurons fire before withdrawal (``on''), others fall silent (``off''); weak stimulation of this area suppresses the reflex. Review: the same circuits both facilitate and inhibit pain transmission, and opioids act through them. & \cite{Fields1983rvm, Fields2004} & measured \\
10 & Expectation of relief dampens pain through the opioid channel & People with acute pain: in those whose pain was reduced by the expectation of relief, blocking the opioid receptors brought the pain back to the level of the others; in the others the blockade did not change the pain. & \cite{Levine1978} & measured \\
11 & The brain chooses the volume of the alarm by the benefit of interrupting & There is no experiment in which the rule for choosing the volume has been measured. & --- & hypothesis \\
12 & Sensitization: after injury the gate output grows and the threshold falls, partly because of the spinal cord; it persists after the damaging stimulus has stopped & Rat: after skin injury the threshold of the flexion reflex falls and the response grows; electrophysiological analysis shows that part of the growth arises in the spinal cord itself. A damaging stimulus causes long-lasting changes in sensitivity that outlast the stimulus. The chapter gives no exact decay time. & \cite{Woolf1983} & measured \\
13 & Sensitization~\eqref{eq:sensitization} is positive feedback; with a strong loop the alarm can latch after healing & Follows from the model~\eqref{eq:sensitization} (an exercise at the end of the chapter). There is no experiment showing that persistent pain after healing is a latched loop of exactly this kind. & derived in the chapter & hypothesis \\
14 & Pain is a negative reward, and learning from it follows the temporal-difference error & Scanner, humans, chains of pain predictors: activity of the ventral striatum and the anterior insula matches temporal-difference learning signals. Monkey: some \nt{DOPA} neurons are excited by an aversive air puff and its predictor, others are inhibited. & \cite{Seymour2004, Matsumoto2009} & measured \\
15 & Pain interrupts the current work & Humans, an electrical pain stimulus and auditory probes: the response to a probe $100\ms$ after pain onset is markedly slowed, and after $1.5\,\text{s}$ there is no longer a difference from control. A review combines such experiments into a model of attentional capture. & \cite{Crombez1996disrupt, EcclestonCrombez1999} & measured \\
16 & Withdrawal is wired in the spinal cord & Rat: neurons of the deep layers of the dorsal horn reproduce the spatial pattern of the withdrawal reflex of individual muscles; they cannot be excited antidromically from the cervical cord, so they are spinal interneurons. & \cite{Schouenborg1995wr} & measured \\
\bottomrule
\end{longtable}}

\section{Chapter~\ref{sec:motorlearning}. Motor Learning}

The least-mean-squares rule and the advantage of a separate error wire are theorems; the structure of the error-wire circuit, weakening on coincidence, tuning of the trace to the delay, the aftereffect, and spontaneous recovery of a skill have been measured; that the whole circuit works as supervised learning and builds an internal model is the leading hypothesis; the numbers of the two-process model are a calculation with the book's model.

{\footnotesize
\setlength{\tabcolsep}{3pt}\renewcommand{\arraystretch}{1.15}
\begin{longtable}{>{\raggedright}p{0.5cm}>{\raggedright}p{4.0cm}>{\raggedright}p{5.6cm}>{\raggedright}p{2.3cm}>{\raggedright\arraybackslash}p{1.7cm}}
\toprule
No. & Claim & Experiment and what was measured & Source & Status \\
\midrule\endfirsthead
\toprule
No. & Claim & Experiment and what was measured & Source & Status \\
\midrule\endhead
1 & The rule~\eqref{eq:lms} on average follows the gradient of the squared error & A result of adaptive-filter theory: a correction proportional to the input and the output error is stochastic gradient descent on the mean squared error. & \cite{WidrowHoff1960} & theory \\
2 & With an error wire for each output, the speed does not depend on the number of outputs; with one shared number, it falls with the number of noisy neurons (Proposition~\ref{prop:teachers}) & Learning curves of a linear network: learning by random perturbation of the neurons is slower than the exact gradient by a factor equal to the number of perturbed neurons. & \cite{Werfel2005} & theory \\
3 & The error-wire circuit works as supervised learning (Proposition~\ref{prop:errorwire}) & Theories derived from the structure of the circuit. The experiments of rows 7--10 agree with them; that the whole circuit works exactly this way has not been proved. & \cite{Marr1969, Albus1971} & hypothesis \\
4 & The expander: about $7\cdot10^{10}$ small \nt{GLUT} neurons, more than in all the rest of the brain & Counting nuclei in a suspension of dissolved tissue: the human cerebellum has $69\cdot10^9$ neurons out of $86\cdot10^9$ in the whole brain. Stereology: $101\cdot10^9$ small neurons in the cerebellum of elderly men. & \cite{Azevedo2009, Andersen1992cereb} & measured \\
5 & Raising the dimension of the input makes the desired dependence linear & The function-counting theorem: a weighted sum of $n$ inputs with a threshold realizes an arbitrary labeling of about $2n$ vectors in general position. That the cerebellum uses exactly this is part of the hypothesis of row~3. & \cite{Cover1965} & theory \\
6 & About $3\cdot10^7$ output \nt{GABA} neurons, each with on the order of $10^5$ inputs & Stereology of the human cerebellum: $30.5\cdot10^6$ output neurons. Electron microscopy, rat: about $1.75\cdot10^5$ inputs from the expander per output neuron. The inhibitory transmitter of the output neurons is from a review. & \cite{Andersen1992cereb, NapperHarvey1988, Ito2001} & measured \\
7 & Exactly one error wire per output neuron, about once per second, each time a powerful burst & A review of recordings: each output neuron receives one climbing \nt{GLUT} input, which evokes a complex spike at a rate of about $1$~Hz. & \cite{Ito2001} & measured \\
8 & The probability of the burst depends on the direction of the movement error & Monkey, oculomotor cerebellum, saccade adaptation: the direction of the visual error sets the probability of a complex spike, and its size sets the timing; the burst changes the simple spikes on the next trial and shifts the movement along the cell's preferred error. & \cite{Herzfeld2018} & measured \\
9 & Coincidence of the error burst with an active input weakens that input ($-\eta e_i x_j$) & A review of experiments: joint stimulation of an input from the expander and the error wire causes long-term weakening of that input; stimulation of the input alone does not. & \cite{Ito2001} & measured \\
10 & The trace length is tuned to the error delay: with a $100\ms$ delay, the input that was active $100\ms$ earlier is weakened most & Slices and living mice: in the region responsible for oculomotor learning, weakening is narrowly tuned to an interval of about $120\ms$ between input and error, the time the error signal takes in this task; in another region different cells are tuned to different intervals. & \cite{Suvrathan2016} & measured \\
11 & The circuit is an adaptive filter~\eqref{eq:adaptivefilter} that learns an internal model of the body & A model derived from the structure of the circuit. There is no direct experiment in which the learned filter has been measured as the transfer function of the body. & \cite{Fujita1982} & hypothesis \\
12 & Prediction subtracts the consequences of one's own movements, so one cannot tickle oneself & Humans: one's own touch feels less ticklish than the same touch by someone else; in the scanner the somatosensory cortex responds more weakly to one's own touch, and the cerebellum is less active when the movement touches the skin. The explanation by subtraction of a prediction is a hypothesis. & \cite{Blakemore1998} & hypothesis \\
13 & Under an external force the miss shrinks, and after the force is removed the arm misses to the other side & People reach with a robot handle in a force field: the trajectories return to straight lines; after the field is removed, the aftereffect is a near mirror image of the initial distortions, and it transfers to untrained parts of the workspace. & \cite{ShadmehrMussaIvaldi1994} & measured \\
14 & Two processes learn~\eqref{eq:tworate}; after a reverse perturbation the skill returns by itself & Humans, a force field, then a short reverse field and trials with the error clamped: the correction returns by itself toward the old one; a two-process model reproduces this, a one-process model does not. & \cite{Smith2006} & measured \\
15 & The numbers of Fig.~\ref{fig:tworate}: $0.84$ (slow $0.63$) after $200$ movements; $6$ reverse ones; fast $-0.53$, slow $0.46$; rebound to $0.37$ within $40$ movements & Calculation with \texttt{models/\allowbreak motor\_\allowbreak learning.py}, $A_f=0.92$, $B_f=0.1$, $A_s=0.996$, $B_s=0.02$: $0.840$ ($0.634$ and $0.205$), $6$ movements, $-0.531$ and $0.457$, a peak of $0.371$ at the $40$th movement. & \texttt{models/\allowbreak motor\_\allowbreak learning.py} & model \\
16 & Two processes are needed because the body changes at different rates & Theory: optimal estimation with disturbances of several lifetimes yields several filters with different time constants and explains spontaneous recovery and faster relearning. & \cite{Kording2007} & hypothesis \\
\bottomrule
\end{longtable}}

\section{Chapter~\ref{sec:chemoutput}. The chemical output}

The two wires to the heart and their different speeds, negative feedback in hormone cascades, the pulse-frequency code, and the daily oscillator and its locking to light are measured directly; the bound $K<8$ is a theorem of control theory derived in the chapter; pulsations generated by the cascade feedback itself are a hypothesis with strong experimental support.

{\footnotesize
\setlength{\tabcolsep}{3pt}\renewcommand{\arraystretch}{1.15}
\begin{longtable}{>{\raggedright}p{0.5cm}>{\raggedright}p{4.0cm}>{\raggedright}p{5.6cm}>{\raggedright}p{2.3cm}>{\raggedright\arraybackslash}p{1.7cm}}
\toprule
No. & Claim & Experiment and what was measured & Source & Status \\
\midrule\endfirsthead
\toprule
No. & Claim & Experiment and what was measured & Source & Status \\
\midrule\endhead
1 & The heart's pacemaker beats by itself about $100$ times a minute; at rest the brake is on, and the rate is usually about $60$--$70$ & Humans, complete block of both wires to the heart: the intrinsic rate is above the resting rate and declines with age, about $100$ beats per minute in adults. So the brake dominates at rest. The resting rate of $60$--$70$ is a typical value, not cited separately. & \cite{JoseCollison1970ihr} & measured \\
2 & The \nt{NORA} gas and the \nt{ACET} brake are low-pass filters, the brake is faster~\eqref{eq:gasbrake} & Dog, frequency-modulated stimulation of the vagal or sympathetic cardiac nerve and the transfer function to atrial rate: both pathways are low-pass filters; the sympathetic one has a much lower corner frequency plus a pure delay of about $1.7\,\text{s}$. The characteristics depend on the mean tone, so the linear formula~\eqref{eq:gasbrake} is an approximation. & \cite{Berger1989} & measured \\
3 & The brake is fast because \nt{ACET} opens the potassium channel without a second messenger, while \nt{NORA} acts through a chain of reactions & Membrane patches of atrial cells: the potassium channel opened by acetylcholine is opened by the $\beta\gamma$ subunits of the receptor-coupled G protein, with no second messenger inside the cell. The \nt{NORA} pathway via cAMP is not cited here. & \cite{Logothetis1987gbg} & measured \\
4 & Blood circulates through the body in about a minute; \nt{ADRE} from the bloodstream reaches every organ with a receptor & A general order-of-magnitude estimate; stated as such in the chapter, no source cited. & --- & estimate \\
5 & A hormone cascade is enclosed in negative feedback: the final hormone inhibits the first stages & Review of experiments: adrenocortical hormones suppress ACTH release by the pituitary and releasing-hormone release by the hypothalamus in three time domains: fast, intermediate, and slow. & \cite{KellerWood1984fb} & measured \\
6 & Three identical stages are stable for $K<8$~\eqref{eq:cascadestable}; at the boundary the period is $2\pi\tau/\sqrt3\approx3.6\tau$ & Derived in the chapter: at the frequency $\sqrt3/\tau$ three stages give a $180^\circ$ phase shift and an attenuation by $8$. & derived in chapter & theory \\
7 & The level error is $1/(1+K)$, so such a regulator cannot do better than about $10\,\%$ (Proposition~\ref{prop:cascade}) & A consequence of row~6 for proportional feedback. The real axis has feedback at several stages and in several time domains (row~5), so the bound applies to model~\eqref{eq:cascade} and is not measured. & derived in chapter & theory \\
8 & At $K=4$ and $7$ the level settles; at $K=12$, steady pulsations with period $3.7$--$3.9\,\tau$ (Fig.~\ref{fig:cascade}) & Computed with \texttt{models/\allowbreak chem\_\allowbreak models.py}: at $K=12$ successive periods are $3.9$, $3.7$, $3.8$, $3.7\,\tau$; at $K=4$ the swing at the end of the run is $0.003$. & \texttt{models/\allowbreak chem\_\allowbreak models.py} & model \\
9 & The stress hormone is released in pulses about once an hour; the pulses are generated by the feedback itself, without a separate oscillator & Rats, constant infusion of releasing hormone: ACTH and corticosterone oscillate at the physiological, nearly hourly frequency, so no rhythmic input from the hypothalamus is needed for the pulses. A pituitary--adrenal model with delayed feedback produces such oscillations. & \cite{Walker2012glc, Walker2010} & hypothesis \\
10 & The receiver reads the pulse frequency, not the level & Monkeys with no release of their own gonadotropin-releasing hormone: constant infusion does not restore the release of pituitary hormones, hourly pulses do; switching back to constant infusion silences the response again. Receptor fatigue as a high-pass filter is an interpretation. & \cite{Belchetz1978} & measured \\
11 & Different pulse frequencies act differently on different cells of one gland & The same monkeys: raising the rate to $2$--$5$ pulses per hour lowers both pituitary hormones; lowering it to one per $3$~hours lowers one (LH) and raises the other (FSH), so their ratio is set by the rate. & \cite{Wildt1981gnrh} & measured \\
12 & The daily rhythm is generated by intracellular negative feedback with a delay of several hours & Review: clock-gene proteins suppress the transcription of their own genes with a delay; the transcription--translation loop gives a period of about a day. & \cite{TakahashiJS2017} & measured \\
13 & The master oscillator is a group of tens of thousands of neurons that synchronizes the rest of the body & Review: the suprachiasmatic nucleus is the master daily oscillator; its individual neurons in culture are autonomous oscillators, and the network synchronizes them; in the mouse, about $10^4$ neurons on each side. & \cite{Welsh2010scn} & measured \\
14 & The intrinsic period in humans is $24.18$~hours & Humans in dim light on a schedule decoupled from the day: the period of the melatonin, temperature, and cortisol rhythms averages $24.18$~hours in young and older people, with a narrow spread. & \cite{Czeisler1999} & measured \\
15 & Light adjusts the oscillator like a phase-locked loop~\eqref{eq:pll}; a shift of about $11$~min per day is needed & Humans, a single $6.7$-hour bright-light exposure at different times of day: a type~1 phase response curve with a $5$-hour range, delay before the body-temperature minimum and advance after it. Equation~\eqref{eq:pll} is the standard model; $11$~min $=0.18$~hours is arithmetic. & \cite{Khalsa2003prc} & measured \\
16 & Without light there is no locking, and the rhythm drifts to its own period & Totally blind people with no light perception living in ordinary society: in about half of them the melatonin and cortisol rhythms run at their own period, which does not match the day. & \cite{Sack1992blind} & measured \\
\bottomrule
\end{longtable}}

\section{Chapter~\ref{sec:debugging}. Debugging the machine}

What an electrode hears, the low dimensionality of activity, the performance of interfaces on rigid arrays, the joint learning of brain and decoder, and the visual spots produced by stimulation are measured in peer-reviewed experiments; the data-rate estimates are the chapter's arithmetic; for the Neuralink device implanted in humans, as far as we could check, data have been published mostly by the company itself.

{\footnotesize
\setlength{\tabcolsep}{3pt}\renewcommand{\arraystretch}{1.15}
\begin{longtable}{>{\raggedright}p{0.5cm}>{\raggedright}p{4.0cm}>{\raggedright}p{5.6cm}>{\raggedright}p{2.3cm}>{\raggedright\arraybackslash}p{1.7cm}}
\toprule
No. & Claim & Experiment and what was measured & Source & Status \\
\midrule\endfirsthead
\toprule
No. & Claim & Experiment and what was measured & Source & Status \\
\midrule\endhead
1 & An electrode hears the spikes of neurons within a radius of about a hundred micrometers, usually one to several; they are told apart by shape & Rat, area CA1, simultaneous intracellular and extracellular recording: the extracellular spike shape follows the width and amplitude of the intracellular one. Estimate: a tetrode could separate $60$--$100$ neurons, but in practice about six are isolated. & \cite{Henze2000extra} & measured \\
2 & The brain has $8.6\cdot10^{10}$ neurons; a probe hears about one hundred-millionth & Counting nuclei in a suspension of dissolved tissue: $86\cdot10^9$ neurons in the human brain. The fraction is the chapter's arithmetic. & \cite{Azevedo2009} & measured \\
3 & The activity of hundreds of motor-cortex neurons fits into one or two dozen shared variables & Review of recordings in monkey motor cortex: population activity is described by a few latent variables shared by many neurons. & \cite{Gallego2017} & measured \\
4 & Neighbors' noise is correlated, and a hundred neurons carry almost as much as a thousand (Proposition~\ref{prop:pooling}) & Monkey visual cortex: the noise correlation coefficient of neighboring neurons is about $0.1$, and the gain from averaging saturates at a hundred neurons. Theory of information-limiting correlations. & \cite{Zohary1994, MorenoBote2014} & measured \\
5 & Raw stream of $2\cdot10^8$~bit/s versus $8\cdot10^4$~bit/s in spikes~\eqref{eq:probedata} & The chapter's arithmetic from the capacity of the spike stream~\eqref{eq:spikeentropy}. The digitization parameters are real: the Neuralink chip samples at $19.3$~kHz, $10$~bits per channel. & derived in chapter; \cite{Rieke1997, Musk2019} & theory \\
6 & The first Neuralink system: chips amplify and digitize, the raw stream goes over a cable, spikes are detected by software outside; on-chip detection, a sealed housing, radio, and wireless power belong to the device for humans, according to the company & Company paper: the full raw stream over a USB-C cable, spikes detected by real-time software off the chip; on-board compression and wireless power and transmission are named as a plan for clinical devices. For the device implanted in humans, only company reports. That on-chip spike detection is necessary is the chapter's conclusion from~\eqref{eq:probedata}. & \cite{Musk2019}; Neuralink reports, no peer-reviewed paper & measured (company paper); for humans, company report \\
7 & Up to $3072$ electrodes on $96$ threads of $32$; a robot inserts the threads, avoiding vessels & Company paper: $3072$ electrodes on $96$ threads of $32$; the robot inserts $6$ threads ($192$ electrodes) per minute; up to $70\,\%$ of electrodes in rats with a chronically implanted array record spikes. & \cite{Musk2019} & measured (company paper) \\
8 & In humans, about a thousand channels on $64$ threads; some threads shifted; cursor at about $10$~bit/s & Company reports only. A PubMed search found no peer-reviewed paper with human results; this agrees with the caveat in the chapter text. & Neuralink reports; no peer-reviewed paper & measured (company report) \\
9 & An interface on an array of a hundred electrodes controls a cursor & A person with paralysis, $96$ electrodes in primary motor cortex: the intention to move the arm modulates spikes three years after the injury; the decoder gives a ``neural cursor'' (email, television), control of a prosthetic hand and a robotic arm. & \cite{Hochberg2006} & measured \\
10 & Writing with the ``mental hand'': about $90$ characters per minute, less than $8$~bit/s & A person with a paralyzed hand, attempted handwriting, recurrent decoder: $90$ characters per minute at $94.1\,\%$ accuracy in real time, above $99\,\%$ after autocorrection. The estimate in bits is the chapter's arithmetic. & \cite{Willett2021} & measured \\
11 & Attempts to speak are turned into text at about $60$ words per minute & A person who lost intelligible speech, arrays in the cortex: $62$ words per minute; $9.1\,\%$ word errors with a $50$-word vocabulary, $23.8\,\%$ with a $125\,000$-word vocabulary. & \cite{Willett2023} & measured \\
12 & An interface extracts a small fraction of the capacity: tens of bits per second for one neuron, thousands for a hundred (Proposition~\ref{prop:probe}) & The upper bound~\eqref{eq:spikeentropy} is theory; the comparison with rows~8 and 10 is the chapter's conclusion. That the bottleneck is low dimensionality and ignorance of the code rather than the wire has not been tested directly. & \cite{Rieke1997} & theory \\
13 & The brain and the decoder learn from each other & Monkeys control a cursor; the decoder adapts in closed loop while the neurons relearn at the same time: accuracy rises, the skill is retained and suffers less from movements of the monkey's own arm. The advice to separate the learning rates is an engineering rule; the chapter has no separate experiment for it. & \cite{Orsborn2014coad} & measured \\
14 & Current through an electrode excites the surrounding neurons and wires regardless of type & Mouse and cat cortex, two-photon calcium imaging: even a weak current excites sparse neurons scattered over up to millimeters, through direct excitation of axons in a volume tens of micrometers across. So the excitation is nonselective but not solid. & \cite{Histed2009stim} & measured \\
15 & Stimulation of the visual cortex lights spots; up to $16$ spots at once allow recognizing some letters and object edges & A blind person, $96$ electrodes in the visual cortex for $6$ months: single-electrode threshold $66.8\pm36.5$~\textmu A; simultaneous stimulation of groups (up to $16$ electrodes, the stimulator's limit) lowers the threshold and gives distinguishable patterns; some letters and object edges are recognized. & \cite{Fernandez2021} & measured \\
16 & Neuralink's second direction is a device that feeds signals into the visual cortex & Only company reports of development; no peer-reviewed data on its performance were found. & Neuralink reports & hypothesis \\
17 & Concentrations of broadcast transmitters can be measured with a chemical sensor in the tissue & Rat brain slices, fast-scan cyclic voltammetry at a carbon fiber: serotonin release and reuptake measured; the substance spills beyond the synapse. & \cite{Bunin1998} & measured \\
\bottomrule
\end{longtable}}

\section{Chapter~\ref{sec:eetypes}. Neurons as devices}

The operating modes of single cells are measured directly, with current injected through an electrode and voltage recording; the mathematics of the integrator and of the division into classes is classical theory, while the idea that the brain uses mode reconfiguration for computation remains a hypothesis.

{\footnotesize
\setlength{\tabcolsep}{3pt}\renewcommand{\arraystretch}{1.15}
\begin{longtable}{>{\raggedright}p{0.5cm}>{\raggedright}p{4.0cm}>{\raggedright}p{5.6cm}>{\raggedright}p{2.3cm}>{\raggedright\arraybackslash}p{1.7cm}}
\toprule
No. & Claim & Experiment and what was measured & Source & Status \\
\midrule\endfirsthead
\toprule
No. & Claim & Experiment and what was measured & Source & Status \\
\midrule\endhead
1 & Resonator: a slow current opposing voltage changes makes the neuron a band-pass filter peaking at a few to tens of hertz (Section~\ref{sec:resonator}) & In slices, sinusoidal current of smoothly rising frequency was injected into \nt{GLUT} neurons and the impedance $|Z(f)|$ was measured: a peak at $2$--$5$~Hz at $33^\circ$C (about $7$~Hz at $38^\circ$C); the resonance is created by slow potassium and cation currents and amplified by a persistent sodium current. A review describes the same technique for many cell classes & \cite{HuVervaekeStorm2002,HutcheonYarom2000} & measured \\
2 & The slow current behaves like an inductance and, together with the membrane capacitance, forms a tank circuit & The subthreshold response of an axon to current was measured and described by linearized conductance equations; the equivalent circuit contains a ``phenomenological'' inductance whose value depends on voltage & \cite{Mauro1970} & theory \\
3 & The time constant of a group with feedback is $\tau_{\text{eff}}=\tau/(1-w)$~\eqref{eq:perfectint}; with $\tau=0.1$~s and $w=0.995$ it is $20$~s & Derived in the chapter from the linear group equation; proposed in a theoretical paper as the mechanism for holding eye position & derived in chapter; \cite{Seung1996} & theory \\
4 & The eye-position integrator holds its input by the network, not by a property of a single cell & Intracellular recording in fish in the nucleus that holds eye position: after a fast eye turn the neuron's rate changes in a step and holds; a brief current into one cell produces only a transient shift, while the voltage steps persist below threshold as well, sustained by synaptic inputs & \cite{Aksay2001} & measured \\
5 & A leak-free integrator needs constant retuning by learning; when mistuned, it either forgets or runs into saturation & Fish were trained with visual feedback proportional to $\pm$ eye position: holding became unstable or leaky, and the neurons' time constant fell by more than an order of magnitude, to $<1$~s; ordinary visual feedback gradually restored stability & \cite{Major2004} & measured \\
6 & Bistable neuron: a brief input switches on a lasting plateau, inhibition switches it off; the property is enabled by serotonin (Section~\ref{sec:bistable}) & Intracellular recording from the \nt{ACET} neurons that control muscles in the cat: brief synaptic excitation puts the neuron into lasting activity, brief inhibition resets it; in the spinal cat the state appears after a serotonin precursor is given & \cite{Hounsgaard1988} & measured \\
7 & Two classes: integrators (rate rises from zero) and resonators (a finite rate right at threshold) & The classes are derived from bifurcation theory. Experiment: in cortical slices, \nt{GLUT} neurons start firing at arbitrarily low rates (type~1), fast \nt{GABA} neurons jump in at a finite rate (type~2) & \cite{Izhikevich2007,Tateno2004} & measured \\
8 & One neuron is an integrator or a coincidence detector: inhibition shortens the summation window & In hippocampal slices, the window within which excitatory inputs sum to threshold is shorter than $2$~ms; it is shortened by feedforward inhibition arriving right after the excitation; in dendrites, where inhibition is weaker, the window is wider & \cite{Pouille2001} & measured \\
9 & A delay line made of axons (Table~\ref{tab:eetypes}) & In the barn owl, delays were measured in the axons leading to coincidence detectors: different axon lengths give different delays, and the detectors are tuned to the difference in sound arrival times & \cite{Carr1990} & measured \\
10 & A pacemaker in wakefulness and a silent cell in sleep & \nt{SERO} neurons fire almost regularly during the day, slow down in slow-wave sleep, and are nearly silent in REM sleep (more in Chapter~\ref{sec:sleep}) & \cite{JacobsAzmitia1992,McGintyHarper1976} & measured \\
11 & The device is set by ion channels and by the broadcast channels that tune them (Proposition~\ref{prop:eetypes}) & Shown in small networks: modulatory substances change the intrinsic properties of neurons and the strength of synapses, so the same circuit produces different rhythms & \cite{Marder2012} & measured \\
12 & The brain uses mode reconfiguration for computation (Proposition~\ref{prop:eetypes}) & There is no direct experiment in which reconfiguring a cell's mode was needed to solve a task; there are only experiments in which modulators change a circuit's output & \cite{Marder2012} & hypothesis \\
\bottomrule
\end{longtable}}

\section{Chapter~\ref{sec:dendrites}. Number of dendrites}

The structure of the classes and the numbers of inputs are measured anatomically and by electrical recordings; the estimate~\eqref{eq:dendriteload} is simple capacitor physics, while the computational explanation of the number of inputs rests on theory and models.

{\footnotesize
\setlength{\tabcolsep}{3pt}\renewcommand{\arraystretch}{1.15}
\begin{longtable}{>{\raggedright}p{0.5cm}>{\raggedright}p{4.0cm}>{\raggedright}p{5.6cm}>{\raggedright}p{2.3cm}>{\raggedright\arraybackslash}p{1.7cm}}
\toprule
No. & Claim & Experiment and what was measured & Source & Status \\
\midrule\endfirsthead
\toprule
No. & Claim & Experiment and what was measured & Source & Status \\
\midrule\endhead
1 & Specific membrane capacitance $c_m\approx1$~\textmu F/cm$^2$ & Membrane was pulled from a neuron's cell body into a pipette, a voltage step was applied, the total capacitance was found from the decay of the capacitive current and then divided by the area: $0.9$~\textmu F/cm$^2$ in three neuron classes & \cite{Gentet2000} & measured \\
2 & Charging time $t\approx c_mA\Delta V/I$~\eqref{eq:dendriteload}: a large neuron needs dozens of simultaneous inputs, a small one is served by a single large synapse & Estimate from the capacitor charging law; the numbers ($20$ and $300$~pF, $0.1$ and several nA) are orders of magnitude & derived in chapter & theory \\
3 & One huge synapse gives a current of several nanoamperes and drives a small neuron to threshold at once & Simultaneous recording of the terminal and the receiver in a slice: single-synapse current $2$--$13$~nA, every input spike evokes a suprathreshold response, delay about $0.4$~ms at $36^\circ$C; the receiver outputs \nt{GLYC} & \cite{Borst1995,BorstSoria2012} & measured \\
4 & Zero dendrites: in touch and pain sensory neurons the spike travels from the sensor to the spinal cord past the cell body & The structure (an axon splitting in two at the cell body) is established anatomically. That conduction does not require an excitable cell body is shown in a validated model of such a neuron: without an excitable body the spike passes the fork just as reliably & \cite{AmirDevor2003} & theory \\
5 & One dendrite: an olfactory neuron carries one receptor type and transmits one input line & Review of experiments with receptor genes: each olfactory neuron expresses one receptor gene from a large family & \cite{Mombaerts2004} & measured \\
6 & Two dendrites: in sound direction detectors, inputs from the left and right ears land on different dendrites & Anatomy and recordings in mammals collected in a review: inputs from the two ears are separated onto two opposite dendrites & \cite{Grothe2010} & measured \\
7 & Separating inputs onto two dendrites makes the AND stricter & Shown in a model: saturation~\eqref{eq:shunt} on one dendrite keeps the inputs from one side from reaching threshold, and coincidence detection improves. There is no direct experiment in which the input layout was changed & \cite{AgmonSnir1998} & theory \\
8 & About $7\cdot10^{10}$ small \nt{GLUT} neurons in the motor learning circuit, with $\sim4$ inputs each & Cell counting by isotropic fractionation: about $69$ billion neurons in this part of the human brain. Recording of one such cell in a live rat: a touch brings bursts of discrete currents from a few inputs, and the cell fires when they sum & \cite{Azevedo2009,Chadderton2004} & measured \\
9 & A threshold element with $n$ inputs separates at most $P_{\max}\approx2n$ random patterns~\eqref{eq:cover}; for the output neuron of the motor learning circuit the bound is $\sim2\cdot10^5$, and a realistic estimate is $\sim4\cdot10^4$ associations & The bound is a classical result of combinatorial geometry. The realistic estimate is capacity theory with positive weights only and a noise margin, applied to the measured number of inputs of this neuron & \cite{Cover1965,Brunel2004} & theory \\
10 & Mixers with few inputs are needed to expand the input; ``four'' is close to optimal & Theory: the dimensionality of the representation given by a layer of mixers is maximal at roughly the number of inputs observed anatomically; the original expansion hypothesis is in the classical papers & \cite{Marr1969,Albus1971,LitwinKumar2017} & hypothesis \\
11 & Large trees: $10^4$--$10^5$ inputs & Synapse counts from electron micrographs: $\approx1.7\cdot10^5$ synapses on one output \nt{GABA} neuron of the motor learning circuit in the rat; about $30\,000$ excitatory and $1\,700$ inhibitory inputs on a hippocampal \nt{GLUT} neuron & \cite{NapperHarvey1988,Megias2001} & measured \\
12 & The branches of a large tree are separate subcircuits with their own thresholds; the neuron is a small multilayer network & Focal stimulation of two sites on thin dendrites: inputs on one branch sum sigmoidally, on different branches linearly. Graded calcium spikes were found in the dendrites of human cortical neurons that let a single cell solve a linearly nonseparable task. Model: reproducing one neuron requires a deep network & \cite{Polsky2004,Gidon2020,Beniaguev2021} & measured \\
13 & Output dendrites: each branch of a retinal \nt{GABA} neuron is a separate direction detector & Two-photon calcium imaging in individual branches: the signal in a branch is selective for motion from the cell body toward the tip, while the somatic voltage is not & \cite{Euler2002} & measured \\
14 & The number of inputs follows the task (Proposition~\ref{prop:dendrites}) & The structure of the classes is measured (rows 3--13); that the number of inputs is chosen for speed or for classification is shown only by theory and models for individual classes & \cite{LitwinKumar2017,AgmonSnir1998} & hypothesis \\
\bottomrule
\end{longtable}}

\section{Chapter~\ref{sec:sleep}. Sleep}

The sleep modes, the replays, and the shrinking of synapses are measured by neuron recordings, microdialysis, and electron microscopy; the sleep schedule and the slow swings are described by the book's models, while the purposes of sleep are mostly hypotheses.

{\footnotesize
\setlength{\tabcolsep}{3pt}\renewcommand{\arraystretch}{1.15}
\begin{longtable}{>{\raggedright}p{0.5cm}>{\raggedright}p{4.0cm}>{\raggedright}p{5.6cm}>{\raggedright}p{2.3cm}>{\raggedright\arraybackslash}p{1.7cm}}
\toprule
No. & Claim & Experiment and what was measured & Source & Status \\
\midrule\endfirsthead
\toprule
No. & Claim & Experiment and what was measured & Source & Status \\
\midrule\endhead
1 & Cortical neurons are not silent in sleep; the operating mode changes & Long recordings of cortical neurons in rats: in slow-wave sleep the rate stays nonzero, and periods of activity alternate with periods of silence; after long wakefulness the rate is higher in all states, after sleep it is lower & \cite{Vyazovskiy2009} & measured \\
2 & \nt{ACET} is high in the day and in REM sleep, low in slow-wave sleep (Table~\ref{tab:sleepmodes}) & Microdialysis in the cortex of freely moving cats: in wakefulness acetylcholine release is $100$--$175\%$ higher than in slow-wave sleep, and in REM sleep it is about at the level of quiet wakefulness & \cite{Marrosu1995} & measured \\
3 & \nt{NORA} and \nt{SERO} quiet down in slow-wave sleep and are nearly silent in REM sleep & Recordings of single \nt{NORA} neurons in rats and \nt{SERO} neurons in cats across sleep stages: the rate is highest in wakefulness, lower in slow-wave sleep, and nearly zero in REM sleep & \cite{AstonJonesBloom1981,McGintyHarper1976} & measured \\
4 & In REM sleep the muscles are switched off by inhibition through \nt{GABA} and \nt{GLYC} & In rats the activity of the \nt{ACET} neurons of a jaw muscle was measured and blockers were applied directly to them: the paralysis is lifted only if both types of \nt{GABA} receptors and the \nt{GLYC} receptors are all blocked & \cite{BrooksPeever2012} & measured \\
5 & Sleep pressure $S$ is a capacitor with two thresholds~\eqref{eq:sleeppressure}, $\tau_r=18.2$~h, $\tau_d=4.2$~h & The two-process model; the time constants are obtained from how EEG slow-wave power rises with wakefulness and falls during sleep, and the shape of the thresholds from the timing of spontaneous waking & \cite{Borbely1982,Daan1984} & theory \\
6 & The machine settles by itself into $16.1$~h awake and $8.0$~h asleep (Fig.~\ref{fig:twoprocess}) & Computed from~\eqref{eq:sleeppressure} with swinging thresholds, script \texttt{models/sleep\_models.py}: $16.1$~h awake and $7.9$--$8.0$~h asleep & --- & model \\
7 & After $40$~h without sleep $S=0.90$, but sleep lasts only $8.8$~h: the debt is repaid by depth, not by duration & Model: \texttt{models/sleep\_models.py} gives $S=0.90$ and $8.8$~h. Experiment: after $40.5$~h without sleep, in eight people EEG slow-wave power in the first recovery night is significantly above normal & \cite{Borbely1981} & model \\
8 & Physically, $S$ is adenosine & Microdialysis in cats: extracellular adenosine in one of the regions that control wakefulness rises during wakefulness, rises further with sleep deprivation, and falls during sleep. That $S$ is adenosine alone has not been shown & \cite{PorkkaHeiskanen1997,Dunwiddie2001} & hypothesis \\
9 & In slow-wave sleep the cortex swings: activity for a few hundred ms, then silence, less than once a second ($<1$~Hz, under anesthesia most often $0.3$--$0.4$~Hz) & Intracellular recording in anesthetized cats: depolarization for $0.8$--$1.5$~s and long hyperpolarization, rhythm $<1$~Hz, most often $0.3$--$0.4$~Hz; the same alternation of activity and silence is seen in natural sleep (row~1) & \cite{Steriade1993,Vyazovskiy2009} & measured \\
10 & The swings are a group with feedback and fatigue~\eqref{eq:updown}; at $b=5$ the period is $1.1$~s (a model value) and activity takes about $35\%$ of the time; at $b=1$ activity is steady (Fig.~\ref{fig:updown}) & Computation, script \texttt{models/sleep\_models.py}: period $1.11$~s, active time fraction $0.35$ (averaged over a whole number of periods); at $b=1$ the active fraction is $1.0$. The model period is shorter than measured (row~9) & --- & model \\
11 & \nt{ACET} switches off the swings and moves the cortex into steady activity & Stimulation of a nucleus of \nt{ACET} neurons in the cat blocked the slow rhythm in $79\%$ of cortical neurons and moved them into steady activity, mainly because the long hyperpolarizations disappeared. Acetylcholine suppresses the slow potassium currents that create fatigue & \cite{SteriadeAmzica1993,McCormick1992} & measured \\
12 & The $7$--$15$~Hz rhythm bursts are generated by the \nt{GLUT}--\nt{GABA} loop between the cortex and a relay node & In slices of the ferret's visual relay node, the bursts are created by rebound bursts of relay cells in response to inhibition from a neighboring \nt{GABA} nucleus, which they themselves excite & \cite{vonKrosigk1993,SteriadeMcCormickSejnowski1993} & measured \\
13 & The day's replays run fast-forward: about $20$ times faster in the hippocampus, $6$--$7$ times in the cortex & In slow-wave sleep, repeating sequences of hippocampal neurons run compressed in time. After running on a track, the sequences of place neurons were replayed in the same order in short bursts of $\sim100$~ms, compressed about $20$ times; in the cortex the compression is $6$--$7$ times & \cite{Nadasdy1999,LeeWilson2002,Euston2007} & measured \\
14 & Strengthening the alignment of replays with the cortex improves memory (a causal link) & In rats after training, stimulation locked to replays strengthened their coupling to the slow waves and rhythm bursts in the cortex: the next day memory was high, while in controls it was at chance & \cite{Maingret2016} & measured \\
15 & The purpose of replays is interleaved training of the slow memory & The theory of two learning systems and computations on artificial networks: sequential learning corrupts the old, interleaved learning does not. There is no direct experiment in the brain showing that replays are needed for exactly this & \cite{McClelland1995} & hypothesis \\
16 & After sleep synapses shrink in proportion to their size; large ones hardly change & Three-dimensional electron microscopy of $6920$ synapses in mouse cortex: the contact area after sleep is $\sim18\%$ smaller than after wakefulness, the reduction is proportional to size, and large stable synapses do not change & \cite{deVivo2017} & measured \\
17 & The main job of sleep is renormalizing the weights & The synaptic homeostasis hypothesis; rows 1 and 16 agree with it but do not prove that this is the main function & \cite{TononiCirelli2014} & hypothesis \\
18 & REM sleep is a bench test of the world model & The mode (Table~\ref{tab:sleepmodes}) is measured; that the network runs a world model is a theoretical proposal with no experiment & \cite{Hobson2009} & hypothesis \\
19 & Flushing of the brain in sleep: the question is open & One experiment: in sleep the extracellular space is $60\%$ wider and clearance is faster. Another, with a different measurement method: in sleep and under anesthesia clearance is markedly slower. The results contradict each other & \cite{Xie2013,Miao2024} & hypothesis \\
20 & Local sleep: patches of an awake animal's cortex briefly drop into silence & In rats after long wakefulness, the neurons of an individual cortical patch briefly fall silent, as in slow-wave sleep, with a slow wave in the local EEG; at such moments the rat makes more errors in a task & \cite{Vyazovskiy2011} & measured \\
\bottomrule
\end{longtable}}

\section{Chapter~\ref{sec:livingmachine}. A Machine Made of Living Neurons}

The behavior of cultures on electrode arrays is measured in direct experiments with recording and stimulation; the burst mechanism rests on the synaptic-depletion model and on experiments in microcultures, and the claims about dopamine in a dish rest on single experiments, some of which the authors themselves regard only as demonstrations.

{\footnotesize
\setlength{\tabcolsep}{3pt}\renewcommand{\arraystretch}{1.15}
\begin{longtable}{>{\raggedright}p{0.5cm}>{\raggedright}p{4.0cm}>{\raggedright}p{5.6cm}>{\raggedright}p{2.3cm}>{\raggedright\arraybackslash}p{1.7cm}}
\toprule
No. & Claim & Experiment and what was measured & Source & Status \\
\midrule\endfirsthead
\toprule
No. & Claim & Experiment and what was measured & Source & Status \\
\midrule\endhead
1 & A culture on a chip: mostly \nt{GLUT}, a smaller fraction of \nt{GABA} that depends on the preparation; about $6\%$ in hippocampal cultures & A network of thousands of neurons on an electrode array is a standard experiment (row~3). The fraction of \nt{GABA} neurons in hippocampal cultures was counted by GABA staining: about $6\%$ of neurons. For cortical cultures we give no verified number & \cite{Benson1994} & measured \\
2 & Organoids: three-dimensional neural tissue about a millimeter across, grown from human stem cells & Three-dimensional organoids with several brain regions were grown from stem cells, including cortex with progenitor layers and mature neurons & \cite{Lancaster2013} & measured \\
3 & With no input a culture fires in bursts at intervals from a second to minutes, depending on the culture & $58$ cortical cultures with densities from $3\,000$ to $50\,000$ neurons were recorded for five weeks ($963$ half-hour recordings): after two weeks, whole-network bursts dominate in almost all cultures; intervals between bursts range from $1$ to $300$~s and depend strongly on the preparation & \cite{Wagenaar2006} & measured \\
4 & A burst ends by transmitter depletion, interval $T_{\text{burst}}\approx\tau_{\text{rec}}\ln\frac{1}{1-x^*}$~\eqref{eq:burstinterval}; $2$--$4$~s is a model estimate with $\tau_{\text{rec}}=0.8$~s, measured recovery is slower & The formula is derived in the chapter. Model: a network with depressing synapses generates whole-network bursts on its own. Experiment in microcultures of $4$--$30$ neurons: burst duration depends on the time since the previous burst, and depleting the transmitter vesicles abolishes bursts; the authors conclude that a burst is limited by the transmitter store. Store recovery there takes tens of seconds, which with the same formula gives intervals of tens of seconds to minutes & derived in the chapter; \cite{Tsodyks2000,CohenSegal2011,Tsodyks1997} & theory \\
5 & ``A teacher who falls silent'': stimulation is switched off when the desired response appears, and the response is locked in & The culture was stimulated at $0.3$--$1$~Hz until the selected response appeared $50\pm10$~ms after the stimulus, then stimulation was switched off for $5$~min; after several cycles the desired response was evoked right away; learning curves were plotted & \cite{Shahaf2001} & measured \\
6 & A culture plays tennis; significant growth of runs within a session only with full feedback & Details in chapter~\ref{sec:dishbrain} and its supplement: with silence in place of the stimulus the culture is also better by the end of the session than without feedback, but with a smaller effect; learning from session to session is not robust & \cite{Kagan2022} & measured \\
7 & That the culture learns to predict its input is a hypothesis; loud claims of ``sentience'' are disputed & The free-energy principle is a theoretical proposal; there is no direct experiment that distinguishes it from simpler explanations. Thirty researchers pointed out in a response article that behavior change in a loop shows neither intelligence nor sentience & \cite{Friston2010,Balci2023} & hypothesis \\
8 & A culture drives a virtual body toward a target with a chosen stimulus pattern & A program chose training stimulation patterns based on the current performance; within tens of minutes the network reached the target states, and plasticity lasted $80$~min after $2$~h of training. The same stimuli delivered without a link to behavior had no effect & \cite{Bakkum2008} & measured \\
9 & Reservoir: only the linear readout $y=W_{\text{out}}\mathbf r$~\eqref{eq:reservoir} is trained & Theory of reservoir computing: any nonlinear system with fading memory plus a trained linear output & \cite{Maass2002,JaegerHaas2004} & theory \\
10 & An organoid as a reservoir identifies speakers with about $78\%$ accuracy (as reported by the authors); the readout learns from the error, and the organoid's connections change without a teacher & Speech sounds from several speakers were delivered to an organoid as current patterns on an electrode array; the trained readout identified the speaker with about $78\%$ accuracy, as the authors report. They also report that the organoid's functional connectivity changed during training and call this unsupervised learning within the organoid itself & \cite{Cai2023} & measured \\
11 & A million neurons spend about $10^{-4}$~W on spikes~\eqref{eq:dishpower} & Estimate: a spike cost of $10^{-10}$~J from the cortical energy budget~\eqref{eq:energy}, times the number of spikes; there is no direct measurement of a culture's power & derived in the chapter; \cite{Attwell2001} & theory \\
12 & Dopamine on a flash of light: $1$~mM caged dopamine, $240$ repetitions, the number of evoked spikes grew & On a remote organoid platform, dopamine in a light-sensitive cage ($1$~mM) was released by ultraviolet light ($365$~nm, $0.8$~s) whenever the stimulus evoked more spikes; over $240$ steps (about $5$~h) the number of spikes grew overall. The authors write that a single experiment cannot establish a link to dopamine; there are no controls & \cite{Jordan2024} & hypothesis \\
13 & Dopamine from living cells changes the weight of an organic transistor lastingly and reversibly & Dopamine-releasing cells were grown above an organic neuromorphic element; oxidation of dopamine at the electrode changes the channel conductance; long-term weight change and its reset were shown & \cite{Keene2020} & measured \\
14 & Organoids with working \nt{DOPA} neurons exist; their dopamine has not been used as a teacher & Electrically active mature \nt{DOPA} neurons and dopamine production were found in organoids grown from human stem cells. We found no experiments in the literature where this dopamine teaches another network & \cite{Jo2016} & measured \\
15 & An organoid balances a pole: training stimuli from a program beat random ones, without consolidation, via \nt{GLUT} synapses & Mouse cortical organoids in a closed loop with the cart-pole task: in most organoids, signals chosen by reinforcement learning gave better performance than random signals or none; after $45$~min of rest the improvement did not persist; blocking AMPA and NMDA receptors abolished it & \cite{Robbins2026} & measured \\
16 & Without control registers a network on the bench cannot decide on its own what counts as success (proposition~\ref{prop:livingmachine}) & In all experiments of rows 5--15 the training signal is set by the experimenter or a program; there is no experiment in which a culture developed a success criterion on its own. This is an inference from absence, not an experiment & --- & hypothesis \\
\bottomrule
\end{longtable}}

\section{Chapter~\ref{sec:dishbrain}. DishBrain}

The bench design and the game results come from one experimental study and its follow-up; the formulas for noise and for drift toward good configurations are derived in the chapter and checked with the book's model, while the learning mechanism in the dish has not been established.

{\footnotesize
\setlength{\tabcolsep}{3pt}\renewcommand{\arraystretch}{1.15}
\begin{longtable}{>{\raggedright}p{0.5cm}>{\raggedright}p{4.0cm}>{\raggedright}p{5.6cm}>{\raggedright}p{2.3cm}>{\raggedright\arraybackslash}p{1.7cm}}
\toprule
No. & Claim & Experiment and what was measured & Source & Status \\
\midrule\endfirsthead
\toprule
No. & Claim & Experiment and what was measured & Source & Status \\
\midrule\endhead
1 & Bench: about $800\,000$ mouse cortical cells or about $10^6$ human cells per chip; $26\,000$ electrodes on $8$~mm$^2$, recording up to $1024$, stimulation through $8$ & Methods of the study: MaxOne chip, $26\,000$ platinum electrodes on $8$~mm$^2$, recording of $1024$ channels at $20\,000$ samples per second; up to $32$ electrodes can technically be stimulated, $8$ could be stimulated independently; mouse cultures were used from about $10$~days & \cite{Kagan2022} & measured \\
2 & A loop with a $10$~ms clock: spikes from the motor regions move the paddle (fig.~\ref{fig:dishloop}) & Spikes are counted in $10$~ms windows; the delay from a spike to the responding stimulus is about $5$~ms, set by the hardware buffer & \cite{Kagan2022} & measured \\
3 & Input: place code (which of $8$ electrodes) and rate code, $4$--$40$~pulses/s & In the main experiment the stimulation rate rises linearly from $4$~Hz when the ball is at the far wall to $40$~Hz at the paddle; the amplitude of the ball pulses is $75$~mV; the pulses are biphasic and rectangular & \cite{Kagan2022} & measured \\
4 & Output $\Delta p=c\,(n_\uparrow-n_\downarrow)$~\eqref{eq:dishreadout}, count noise per window $1/\sqrt{RT}$ & The two-region difference rule is described in the study (with region weights by activity); the exact formula is not given. The noise estimate follows from Poisson counting and is derived in the chapter & \cite{Kagan2022}; derived in the chapter & theory \\
5 & Teacher: after a hit, a predictable stimulus of $75$~mV, $100$~pulses/s, $0.1$~s; after a miss, an unpredictable one, $150$~mV, $5$~pulses/s, $4$~s at random places and times, then a $4$~s pause & Methods: after a hit, $75$~mV, $100$~Hz, $100$~ms to all $8$ electrodes at once; after a miss, $150$~mV, $5$~Hz to random electrodes at random times for $4$~s, then $4$~s without stimulation. The culture contains no dopamine & \cite{Kagan2022} & measured \\
6 & The network acts so as to make its input predictable & The free-energy principle is a theoretical proposal from which the authors derived the protocol; the experiment agrees with it but does not distinguish it from simpler mechanisms (rows~7--8) & \cite{Friston2010,Kagan2022} & hypothesis \\
7 & If failure shakes the network and success does not, the time in a configuration is $\rho\propto1/P_{\text{miss}}$~\eqref{eq:dishdensity} (proposition~\ref{prop:dishscramble}) & Follows from flux balance in a Markov chain with symmetric kicks; derived in the chapter. There is no direct test on a culture & derived in the chapter & theory \\
8 & The ``scramble on a miss'' model learns: fraction of hits $0.21\to0.38$; with a kick after every rally $\approx0.12$, with no kicks $0.19$ (fig.~\ref{fig:dishmodel}) & Computation, script \texttt{models/dishbrain\_model.py}, $40$ model cultures with $2000$ rallies each: $0.21\to0.38$; $0.13\to0.11$; $0.19\to0.19$ & --- & model \\
9 & Runs of hits grow longer only in living cultures in the full loop; the controls do not learn & $399$ sessions of $20$~min: $9$ mouse cultures ($101$ sessions), $11$ human ($138$), $6$ medium-only chips ($80$), $20$ cultures at rest ($42$), $3$ simulations ($38$). The mean rally length over the last $15$~min exceeds that over the first $5$ only in living cultures; non-neuronal cells (HEK293T) and medium-only chips show no effect & \cite{Kagan2022} & measured \\
10 & Significant growth within a session only with full feedback; with silence in place of the stimulus, play is better by the end of the session than without feedback, but the effect is smaller & $486$ sessions over $3$ days: ``stimulus'' ($n=27$), ``silence'' ($17$), ``no feedback'' ($15$). Growth over time is significant only in the ``stimulus'' condition; at the end of the session ``silence'' outperformed ``no feedback'' with a smaller effect & \cite{Kagan2022} & measured \\
11 & The effect is small; whether the network learns for the long term is an open question & Within a session the improvement is reliable, but, as the authors themselves say, learning across sessions over several days was not robustly observed & \cite{Kagan2022} & measured \\
12 & During play the network approaches a critical regime, and the closer it gets, the better the play & Analysis of activity avalanches in the same cultures: structured input brings the network closer to the critical regime, and play quality correlates with proximity to it; but criticality alone, without information about the consequences of actions, is not enough for learning & \cite{Habibollahi2023} & measured \\
13 & The culture's ``sentience'' has not been shown & A response article by thirty researchers: behavior change in a closed loop proves neither intelligence nor sentience; no experiment shows it & \cite{Balci2023} & hypothesis \\
14 & Proposed fourth control: a predictable stimulus after a miss, of the same duration and energy (section~\ref{sec:dishspec}) & No such experiment has been done; this is the book's proposal & --- & hypothesis \\
\bottomrule
\end{longtable}}

